\documentclass[11pt,a4paper]{article}
\usepackage{fontspec}
\setmainfont{Malgun Gothic}
\setlength{\oddsidemargin}{0pt}
\setlength{\evensidemargin}{0pt}
\setlength{\textwidth}{6.4in}
\setlength{\topmargin}{-0.35in}
\setlength{\textheight}{9.5in}
\XeTeXlinebreaklocale "ko"
\XeTeXlinebreakskip=0pt plus 1pt
\usepackage{amsmath,amssymb}
\usepackage{longtable,array}
\usepackage{xcolor}
\newcommand{\href}[2]{#2\footnote{\texttt{\detokenize{#1}}}}
\setlength{\parindent}{0pt}
\setlength{\parskip}{0.55em}
\emergencystretch=3em
\setcounter{tocdepth}{2}
\renewcommand{\arraystretch}{1.25}
\begin{document}
\begin{titlepage}
\centering
\vspace*{3cm}
{\Large AI금융경제\par}
\vspace{1.2cm}
{\Huge\bfseries 슬라이드별 학습 가이드\par}
\vspace{1.2cm}
{\large 값싼 예측에서 종목 간 기대수익률까지\par}
\vfill
{\large 2026년 2학기\par}
\end{titlepage}
\tableofcontents
\clearpage

\section{01. 값싼 예측과 AI의 경제학}
인공지능을 금융에 적용한다는 말은 먼저 \textbf{어떤 정보를 얻어 어떤 선택을 바꾸는가}를 묻는 일이다. 예측의 값은 정확도 자체가 아니라 그 예측으로 피할 손실과 새로 얻을 편익에 있다. 이 장은 예측의 경제학에서 과업의 보완관계, 가격 조정, 자산가격 예측으로 이어진다. 수학에서 필요한 것은 조건부기대값, 탄력성, 회귀의 기초이며, 각각을 아래에서 다시 정의한다.\par

\subsection{슬라이드 01 · 강좌의 지도}
금융에서 AI는 공시 읽기, 신용심사, 거래, 위험관리처럼 성격이 다른 일을 한다. 이들을 묶는 질문은 \textbf{미래의 불확실한 상태를 더 잘 추정하면 현재의 행동이 어떻게 달라지는가}다. 예를 들어 다음 분기 매출을 맞혀도 매수·매도 기준이나 거래비용을 정하지 않으면 투자전략은 완성되지 않는다.\par

이 장의 순서는 예측의 가치 → 업무 병목 → 가격과 고용의 변화 → 실제 수익률 예측의 검증이다. 뒤의 정규화 회귀와 요인 논쟁은 앞의 경제학적 질문에 대한 통계적 도구로 읽으면 연결된다.\par

\subsection{슬라이드 02 · 왜 금융에서 AI를 배우는가}
금융의 입력은 공시, 뉴스, 가격, 재무제표처럼 풍부하지만 목표인 \textbf{미래 초과수익률}은 매우 시끄럽다. 같은 기술이라도 영수증 분류에서는 높은 정확도가 곧 생산성으로 이어질 수 있지만, 투자에서는 작은 예측 오차가 큰 손실로 번진다.\par

따라서 이 분야는 AI의 장점을 가장 필요로 하면서도 검증 기준이 특히 엄격하다. 독자는 매번 데이터가 언제 알려졌는지, 선택 가능한 행동이 무엇인지, 성과의 비교 대상이 무엇인지 묻는 습관을 갖추어야 한다.\par

\subsection{슬라이드 03 · 주식 매수 질문에 빠진 변수들}
“삼성전자를 지금 사도 되는가”라는 질문은 답을 하나로 정할 수 없다. 투자 기간, 이미 가진 종목, 손실 감내 범위, 세금과 수수료가 빠졌기 때문이다. 같은 전망을 가진 두 사람도 현재 포트폴리오가 다르면 최적 행동이 다르다.\par

판단의 기본형은 상태 $s$가 일어날 확률 $p(s)$와 행동 $a$의 효용 $U(a,s)$를 놓고 $a^*=\arg\max_a\sum_s p(s)U(a,s)$를 선택하는 것이다. AI는 주로 $p(s)$를 바꾸지만, 목표와 제약을 정하는 일은 별개다. 자동 주문 권한까지 연결된다면 잘못된 예측의 비용과 사전 통제가 훨씬 중요해진다.\par

\subsection{슬라이드 04 · AI 충격과 한국 주식시장의 집중}
한국 주가지수는 시가총액 가중이므로 소수 반도체 대형주의 변동이 지수 전체에 크게 반영된다. 미국의 장기금리 상승은 먼 미래의 이익을 현재가치로 환산하는 할인율을 높이고, AI 투자 수익성에 대한 의문은 예상 현금흐름을 낮출 수 있다. 두 힘이 동시에 작용하면 관련 기업의 가격은 더 민감하게 움직인다.\par

AI 기업 사이의 투자·클라우드 구매·반도체 공급 계약은 서로의 매출을 키우지만 현금흐름의 독립성을 보장하지 않는다. 한 고리의 투자 축소가 다른 고리의 매출 전망을 낮추는 \textbf{공통 위험}이 될 수 있다. 단순한 “AI 수혜주” 분류보다 자금의 출처와 최종 고객 수요를 추적해야 한다.\par

\subsection{슬라이드 05 · 값싼 예측의 경제학}
예측은 아직 관찰하지 못한 값을 현재 정보로 추정하는 것이다. 비용이 내려가면 기업은 예전에는 예측 문제로 보지 않았던 업무까지 예측 가능한 조각으로 분해한다. 대출 심사에서는 “상환할까”뿐 아니라 “어떤 추가 서류가 판정을 바꿀까”도 예측 대상이 된다.\par

가격이 떨어진 재화의 사용량이 늘고, 함께 쓰는 보완재의 가치가 바뀐다는 미시경제학을 그대로 적용할 수 있다. 데이터 품질, 손실함수의 설계, 행동할 권한이 대표적인 보완재다. \href{https://www.predictionmachines.ai/}{Prediction Machines}의 관점은 AI를 만능 지능으로 부르기보다 어떤 투입의 가격이 내려갔는지 묻게 한다.\par

\subsection{슬라이드 06 · 세 명제의 연결}
첫째, 예측은 결정을 바꾸어야 가치가 있다. 둘째, 그 결정이 실제 결과로 이어지려면 업무의 다른 단계가 따라와야 한다. 셋째, 예측 비용이 충분히 내려가면 고객의 수요와 관련 노동·데이터의 가격이 조정된다.\par

세 명제는 개인의 선택, 기업의 생산과정, 시장 균형이라는 서로 다른 분석 수준이다. 한 수준에서 관찰한 성능을 다른 수준의 성과로 곧장 옮기면 논리적 도약이 생긴다. 가령 문서 요약 속도가 10배여도 승인위원회의 처리량이 그대로라면 기업의 총 처리량은 거의 안 바뀔 수 있다.\par

\subsection{슬라이드 07 · 의사결정을 바꾸지 않는 신호의 가치는 0}
신호 $z$를 보기 전 최적 기대효용은 $V_0=\max_a E[U(a,s)]$다. 신호를 본 뒤에는 신호마다 행동을 다시 고를 수 있으므로 $V_1=E_z[\max_a E[U(a,s)\mid z]]$다. 기존 행동을 계속 고르는 선택도 가능하므로 $V_1\ge V_0$이고, \textbf{정보의 가치}는 $V_1-V_0$이다.\par

두 신호의 정답률이 달라도 최적 행동이 늘 같다면 행동 가치의 차이는 없다. 예컨대 투자 허용 범위가 너무 좁아 어떤 신호에도 같은 보유량을 유지한다면 더 정밀한 수익률 전망의 즉각적 가치는 0이다. 반대로 드문 파산을 맞히는 신호는 전체 정확도가 낮아도 큰 손실을 피하게 해 가치가 클 수 있다. 이는 \href{https://doi.org/10.1214/aoms/1177729032}{Blackwell의 정보 비교}를 금융 의사결정에 적용한 것이다.\par

\subsection{슬라이드 08 · 지도와 압축}
추상화는 정보를 버리는 일이다. 피카소의 황소 그림은 세부 묘사를 지우면서도 황소를 알아볼 특징을 남긴다. 지하철 노선도는 실제 거리와 방향을 왜곡해 환승 경로를 읽기 쉽게 한다.\par

좋은 압축의 기준은 원본을 얼마나 완벽히 재현하는지가 아니라 \textbf{목표 결정에 필요한 차이를 보존하는가}다. 지하철 지도는 환승에는 유용하지만 도보 시간에는 부정확할 수 있다. 금융 모형도 종목의 모든 정보를 저장할 필요는 없지만, 손익을 좌우하는 꼬리위험이나 발표 시점을 버려서는 안 된다.\par

\subsection{슬라이드 09 · 한 결정에 맞춰 만든 지도}
현실 $x$를 표현 $z=f(x)$로 압축한 다음 예측 $\hat y=g(z)$를 만들고 행동 $a=h(\hat y)$을 고른다고 하자. $f$는 두 현실을 같게 보이게 할 수 있다. 두 현실에서 최적 행동이 같다면 그 손실은 작지만, 하나는 매수하고 다른 하나는 매도해야 한다면 치명적이다.\par

거대언어모형은 학습자료의 패턴을 압축해 답을 생성한다. 금융에서는 문장 요약의 매끄러움보다 의사결정에 필요한 수치, 부정문, 예외 조항, 기준일을 유지하는지가 중요하다. 공시의 “일회성 이익” 한 단어가 사라지면 이익의 지속성 평가가 달라진다.\par

\subsection{슬라이드 10 · 공시에서 글자가 이미지일 때}
문자 검색은 텍스트 층에 존재하는 글자를 찾는다. 실적 발표의 핵심 문단이 그림으로 삽입되면 사람이 보기에는 같은 문장이지만 기계의 단순 검색에서는 사라질 수 있다. 멀티모달 인식은 이미지의 글자를 읽어 이런 누락을 줄인다.\par

그러나 읽기와 평가를 구분해야 한다. 자산 매각 이익을 정확히 추출했더라도 그 이익이 반복 가능한 영업이익인지, 현금 유입인지, 주가에 이미 반영됐는지는 추가 판단이 필요하다. 실무적으로는 이미지 OCR 결과를 원문 페이지와 대조하고, 추출된 숫자의 단위와 회계 기간을 검증한다.\par

\subsection{슬라이드 11 · 빠른 요약과 정확한 요약의 선택}
분석가가 3일 걸려 오류율 1\%인 결과를 내고 AI가 10초 걸려 오류율 5\%인 결과를 낸다고 하자. “어느 쪽이 좋은가”는 오류 하나의 비용, 마감 시각, 사람이 후속 검토할 수 있는지에 달려 있다. 단순히 속도나 정확도 한 숫자로 고르지 않는다.\par

오류의 평균 비용을 $L$, 지연 비용을 $D$라 하면 총 기대비용의 단순형은 $C=p_{\mathrm{err}}L+D$다. 중요한 공시를 잘못 읽어 큰 주문을 넣는 경우에는 $L$이 커서 검토 단계가 가치 있다. 반면 수천 개 문서에서 검토 후보만 찾는 1차 선별에는 빠른 도구가 유리할 수 있다.\par

\subsection{슬라이드 12 · 암달의 법칙과 업무 병목}
전체 처리시간을 1로 정규화하고, 자동화 가능한 부분이 $p$, 그 부분의 속도가 $s$배가 된다고 하자. 새 시간은 자동화하지 못한 $1-p$와 빨라진 $p/s$의 합이다. 따라서 전체 속도 증가는 다음과 같다.\par

\[
S=\frac{1}{(1-p)+p/s}.
\]

문서 초안이 업무시간의 10\%라면 이를 10배 빠르게 해도 $S=1/(0.9+0.01)\approx1.099$, 즉 전체 처리량 증가는 약 9.9\%다. 준법 검토와 최종 승인에 긴 시간이 걸리면 초안 작성의 혁신이 곧바로 전체 혁신이 되지 않는다. $s\to\infty$일 때조차 상한은 $1/(1-p)$다.\par

\subsection{슬라이드 13 · 영상의학 사례가 알려 주는 것}
영상 판독 성능이 높아졌다는 사실만으로 의사의 고용이 사라진다고 예측할 수 없다. 직업은 판독, 환자와의 소통, 다른 전문의와의 협업, 시술, 법적 책임의 묶음이다. AI가 한 과업을 싸게 만들면 전체 검사량이 늘고 남은 과업의 수요가 커질 수 있다.\par

다만 다른 시기나 국가의 보수·고용 변화를 AI 하나의 인과효과로 읽을 수는 없다. 인구 고령화, 의료 수가, 양성 정원, 규제도 움직인다. 이 사례의 교훈은 “자동화 가능한 과업의 비중”과 “검사 수요의 탄력성”을 함께 살피라는 것이다.\par

\subsection{슬라이드 14 · 암달 곡선의 평탄해짐}
$p=0.5$이면 자동화 부분을 무한히 빠르게 해도 전체는 최대 2배다. $p=0.9$라면 상한은 10배다. 속도 증가의 한계효과는 $s$가 커질수록 작아진다. 분모의 $p/s$가 이미 작아진 뒤에는 다른 단계가 시간을 지배하기 때문이다.\par

실무 결정은 가장 느린 도구를 무조건 개선하는 것이 아니다. 전체 시간에서 해당 단계가 차지하는 비중을 측정하고, 개선 후 새 병목이 어디로 이동하는지 계산한다. 거래 시스템에서 뉴스 분류를 빨리 해도 주문 한도 승인이나 결제 통제가 그대로면 총 처리량은 제한된다.\par

\subsection{슬라이드 15 · 과업 생산성과 총량 효과의 차이}
고객상담, 글쓰기, 소프트웨어 개발 같은 개별 과업에서 AI가 시간을 줄인다는 실험은 중요하다. 그러나 임금·고용·GDP는 기업의 도입률, 품질 검토, 수요 증가, 가격 하락, 다른 과업으로의 재배치가 합쳐진 균형 결과다. 개별 과업의 처리속도 증가를 총량 생산성 증가와 동일시할 수 없다.\par

측정에도 시차가 있다. 초기에는 설치·학습·프로세스 재설계 비용이 들고, 나중에야 보완투자가 이루어진다. 반대로 총량 지표에 변화가 아직 작다고 해서 미시적 편익이 없다는 뜻도 아니다. \href{https://doi.org/10.1093/qje/qjae044}{생성형 AI와 상담 생산성 연구}는 과업 수준의 비교에 유용하다.\par

\subsection{슬라이드 16 · 싸진 예측이 재평가하는 것}
예측 서비스의 한계비용이 내려가면 가격도 경쟁을 통해 내려갈 수 있다. 고객은 같은 예산으로 더 많은 예측을 이용하고, 예측과 함께 쓰는 희소한 데이터·판단·집행 권한의 가치는 오를 수 있다. 다만 누구에게 이득이 귀속되는지는 수요와 공급의 탄력성에 달려 있다.\par

노동은 한 덩어리가 아니다. 자동화된 과업을 많이 수행하는 노동과 예측을 해석·책임지는 노동은 보완 또는 대체 관계가 다르다. 따라서 “AI가 임금을 올린다/내린다”보다 어떤 과업의 수요가 늘고 공급이 얼마나 빨리 따라오는지 묻는 편이 정확하다.\par

\subsection{슬라이드 17 · 제번스 역설의 그림}
에너지 효율이 개선되면 같은 이동 거리의 연료비가 내려간다. 비용이 낮아져 주행거리가 늘 수 있으므로 차량 한 대의 효율과 사회 전체의 연료 사용량은 반대 방향으로 움직일 수 있다. 그림의 혼잡한 도로는 \textbf{단위당 투입 감소와 총투입 증가가 양립}하는 경우를 보여 준다.\par

AI에도 같은 질문을 던진다. 공시 분석 한 건의 비용이 1/10이 되어 분석 건수가 20배가 되면 총 분석 투입은 오히려 2배다. 하지만 이런 반등은 필연이 아니다. 예측 수요가 포화되어 있거나 규제상 주문 횟수를 늘릴 수 없다면 총사용량이 줄 수도 있다.\par

\subsection{슬라이드 18 · 탄력성으로 보는 반등}
단위 서비스 비용이 효율 개선으로 $d\ln P<0$만큼 낮아졌다고 하자. 수요의 가격탄력성을 $\varepsilon=-d\ln Q/d\ln P>0$로 정의하면 서비스 사용량은 $d\ln Q=-\varepsilon d\ln P$만큼 늘어난다. 한 서비스에 필요한 자원이 가격 하락률만큼 줄어든다면 총자원 $E$의 변화는 $d\ln E\approx(1-\varepsilon)d\ln P$다.\par

$\varepsilon>1$이면 $d\ln P<0$에 대해 $d\ln E>0$가 되어 반등이 효율 절감을 넘어선다. 의료영상 검사가 저렴해지면 검사 대상을 넓힐 수 있지만, 보험 규칙과 진료 지침이 수요를 제한할 수도 있다. 탄력성은 사례마다 추정해야 한다.\par

\subsection{슬라이드 19 · 기술과 직업의 역사}
증기기관, ATM, 스프레드시트는 모두 한 과업의 비용을 낮췄다. 은행 창구에서는 현금 인출 업무가 줄어도 지점 운영과 고객 상담의 상대적 중요성이 커질 수 있었다. 회계에서는 계산 자동화와 함께 감사·해석 업무가 늘었다.\par

같은 역사로 모든 직업의 미래를 단정하면 안 된다. 기술 도입 시기에는 지점 수, 규제, 소득 증가가 함께 움직인다. 비교할 것은 직업명보다 과업의 구성, 서비스 수요의 가격탄력성, 남은 과업에 대한 진입 장벽이다.\par

\subsection{슬라이드 20 · 약한 고리와 GDP}
경제 전체의 산출을 서로 보완적인 여러 부문의 함수로 보면 작은 부문의 생산성을 무한히 높여도 전체 산출은 다른 부문에 묶인다. 암달의 법칙이 시간의 합으로 이를 보여 준다면, 보완적인 생산함수는 \textbf{가치의 합과 대체 가능성}까지 고려한다.\par

대체탄력성 $\sigma<1$은 한 부문의 산출을 다른 부문의 산출로 바꾸기 어렵다는 뜻이다. 법률 검토가 빨라져도 법원 절차가 일정하면 전체 사건 처리량의 증가는 작다. GDP 효과를 추정할 때는 자동화 대상의 부가가치 비중과 대체탄력성을 함께 정해야 한다. \href{https://web.stanford.edu/~chadj/JonesTonetti_Automation.pdf}{Jones와 Tonetti의 생산성 논의}는 이 제약을 수량화한다.\par

\subsection{슬라이드 21 · 분석 가격이 내려갈 때 희소해지는 것}
공시의 숫자와 문장을 찾아 정리하는 비용이 크게 낮아지면 예전에는 전문 서비스 구매자만 하던 검색을 더 많은 사람이 할 수 있다. 다만 일반 AI 구독료와 전문 데이터 서비스의 가격을 단순히 나누어 \textbf{동일한 품질의 가격비}로 해석해서는 안 된다. 데이터 범위, 오류 책임, 라이선스, 감사 가능성이 다르다.\par

예측이 보편화될수록 차별적 데이터, 적절한 질문, 거래 실행, 사후 검증이 중요해진다. 동일한 공개 공시에서 같은 숫자를 추출하는 능력만으로 장기 초과수익을 얻기는 어렵다. 정보가 가격에 얼마나 빨리 반영되는지도 함께 봐야 한다.\par

\subsection{슬라이드 22 · 내 일의 위험을 평가하는 세 질문}
첫째, 일의 모든 과업을 자동화할 수 있는가. 둘째, 가격이 내려가면 그 서비스의 수요가 충분히 늘어나는가. 셋째, 남은 보완 과업의 공급이 제한되어 있는가. 세 질문을 함께 보아야 자동화와 임금의 방향을 추론할 수 있다.\par

예를 들어 신용분석에서 숫자 입력은 자동화되어도 고객의 특수 상황 확인과 승인 책임은 남을 수 있다. 하지만 이 남은 업무를 누구나 빠르게 배울 수 있다면 임금 프리미엄은 크지 않을 수 있다. 반대로 규제 자격이나 희소한 관계망이 필요하면 가치가 높아질 여지가 있다.\par

\subsection{슬라이드 23 · 데이터 모형에서 예측 모형으로}
전통 통계는 자료가 어떤 확률모형에서 생겼는지 먼저 가정하고 모수를 추정한다. 알고리즘적 접근은 생성 과정을 완전히 알기 어렵다고 보고, 새로운 자료에서 얼마나 잘 맞히는지 우선 평가한다. 두 관점은 적이 아니다. 예측은 행동에 유용하고, 구조적 해석은 정책 반사실과 원인 분석에 필요하다.\par

금융에서는 예측 성공이 특히 어려워서 두 관점 모두 검증이 필요하다. 이 장의 나머지는 요인모형의 이론적 절제와 머신러닝의 유연성을 같은 표본외 기준으로 비교한다.\par

\subsection{슬라이드 24 · 브라이먼의 두 문화}
\textbf{데이터 모형 문화}는 예컨대 $y=X\beta+\varepsilon$을 가정해 $\beta$의 의미와 유의성을 해석한다. \textbf{알고리즘 문화}는 함수 $f(X)$를 유연하게 고르고 미사용 자료의 예측오차를 줄인다. 전자는 설명 가능한 모수에 강하고 후자는 복잡한 비선형 패턴에 강할 수 있다.\par

여러 모형이 거의 같은 예측력을 갖는 라쇼몽 효과가 있으면 최고 점수 하나를 ‘진실’로 해석하기 어렵다. 단순성과 정확도의 상충도 자료와 목적에 따라 달라진다. \href{https://doi.org/10.1214/ss/1009213726}{Breiman의 원문}을 읽을 때 모델의 목적이 설명인지 결정인지 먼저 구분하자.\par

\subsection{슬라이드 25 · 두 문화의 원문을 읽는 법}
브라이먼의 비판은 확률모형이 언제나 틀리다는 주장이 아니다. 현실이 가정한 형태와 다를 때 모형 내부의 깔끔한 표준오차만 믿으면 새 자료에서 실패할 수 있다는 지적이다. 모형의 타당성을 평가하려면 잔차, 안정성, 표본외 성능을 함께 봐야 한다.\par

금융 연구에서 회귀계수가 유의한 것은 특정 가정 아래 표본 내에서 0과 다르다는 뜻이다. 내년에도 같은 부호와 크기를 유지하거나, 거래비용을 넘는 이익을 준다는 결론은 별도의 검증이 필요하다.\par

\subsection{슬라이드 26 · 금융이 이론 중심 접근을 고수한 이유}
고객과 감독당국은 왜 신용을 거절했는지, 왜 특정 위험을 보유하는지 설명을 요구한다. 또한 금융 수익률의 신호는 약하므로 자유도가 높은 모형은 과거 우연을 외우기 쉽다. CAPM 같은 이론은 어떤 위험에 보상이 주어져야 하는지 검증 가능한 제약을 제시한다.\par

또 한 가지는 정책과 시장의 반응이다. 예측 규칙이 널리 쓰이면 가격이 바뀌어 규칙의 성능 자체가 달라진다. 이는 루커스 비판과 닮았다. 과거의 조건부 상관을 그대로 미래의 구조관계로 취급하지 않는 이유다.\par

\subsection{슬라이드 27 · 요인모형의 언어}
종목 $i$의 초과수익률 기대값을 $E[r_i]=\beta_i^\top\lambda$로 쓴다. $\beta_i$는 시장·가치·수익성 같은 공통 충격에 대한 노출이고 $\lambda$는 그 충격 한 단위에 대해 요구되는 보상이다. CAPM에서는 시장 요인 하나만 두므로 $E[r_i-r_f]=\beta_{iM}E[r_M-r_f]$가 된다.\par

이 식은 \textbf{위험 노출의 가격}과 \textbf{단순 예측 특성}을 구분한다. 장부가치/시가총액이 높은 종목의 수익률이 높아도 그것이 어떤 위험의 보상인지, 지속적 가격오류인지 바로 결정되지 않는다. 적은 수의 요인과 엄격한 경제적 해석은 장점이지만 놓치는 상호작용이 있을 수 있다.\par

\subsection{슬라이드 28 · 요인 동물원과 다중검정}
수백 개 특성을 하나씩 검정하면 우연히 높은 $t$값이 나오는 특성이 반드시 생긴다. 참효과가 모두 0이고 독립 검정 300개를 각각 5\% 수준에서 하면 기대 거짓 발견 수는 $300\times0.05=15$개다. 실제 논문 선택·출판 과정은 이보다 복잡하지만 기본 문제는 같다.\par

따라서 “한 논문에서 $|t|>2$”는 충분하지 않다. 독립 표본, 발표 뒤 성과, 거래비용, 관련 변수 통제, 검정 횟수를 확인해야 한다. \href{https://doi.org/10.1093/rfs/hhv059}{Harvey·Liu·Zhu의 다중검정 연구}는 더 높은 발견 기준의 필요성을 제기한다.\par

\subsection{슬라이드 29 · 머신러닝의 동일 경기장}
Gu·Kelly·Xiu는 다수 미국 주식과 기업 특성을 사용해 선형 회귀, 벌점 회귀, 차원 축소, 나무, 신경망을 비교했다. 핵심은 같은 정보 시점과 같은 미래 기간을 기준으로 성능을 채점하는 것이다. 모형이 복잡할수록 과거 적합도가 높다는 사실만으로 승자를 정하면 불공정하다.\par

특성과 거시 상태의 상호작용은 “같은 기업 특성도 경기 상태에 따라 의미가 달라진다”는 가설을 표현한다. 예를 들어 레버리지의 가격 효과가 금리 상승기에 커질 수 있다. \href{https://doi.org/10.1093/rfs/hhaa009}{연구 원문}은 이러한 비선형 상호작용이 예측 개선의 중요한 원천이라고 보고한다.\par

\subsection{슬라이드 30 · 작은 표본외 결정계수의 의미}
표본외 결정계수는 $R^2_{\rm OOS}=1-\mathrm{SSE}_{\rm model}/\mathrm{SSE}_{\rm benchmark}$이다. 월별 개별 주식 수익률은 예측 불가능한 변동이 크므로 0.4\% 같은 숫자는 작아 보여도 실제 예측 가능한 성분의 차이를 나타낼 수 있다. 단, 통계적 개선과 거래 가능한 이익은 다르다.\par

모형의 우수성을 읽을 때 기준 예측이 0인지 과거 평균인지 확인해야 한다. 분모가 달라지면 같은 예측에도 $R^2$가 달라진다. 공매도 제약, 회전율, 거래비용을 빼고 계산한 포트폴리오 성과도 별도로 검토한다.\par

\subsection{슬라이드 31 · 공개된 예측 신호의 이익 감소}
한 뉴스 신호가 널리 알려지면 투자자는 가격에 그 정보를 더 빨리 반영한다. 신호가 가리킨 종목을 모두가 먼저 사면 나중에 얻을 초과수익은 줄어든다. 거래 규모가 커질수록 호가 충격과 매매비용도 늘어난다.\par

그래서 발표 전후 샤프비율의 하락은 가능한 메커니즘을 보여 주지만, 표본 기간·거래비용·전략 변경 여부를 확인하지 않고 ‘완전 소멸’로 단정할 수 없다. 샤프비율은 평균 초과수익을 표준편차로 나눈 값이며, 비용을 빼기 전 값과 실제 투자자의 값은 다르다. \href{https://doi.org/10.1111/jofi.12365}{McLean·Pontiff}는 발견과 출판 뒤의 감소를 체계적으로 조사한다.\par

\subsection{슬라이드 32 · 계산 규모와 사전 지식}
서튼의 “쓴 교훈”은 사람이 세세한 규칙을 집어넣는 방법보다 계산과 학습을 넓게 활용하는 일반적 방법이 장기적으로 우세했던 사례를 정리한다. 금융의 손으로 만든 단어 사전이나 고정된 요인 조합도 이 관점에서 재검토할 수 있다.\par

그렇다고 경제 이론을 버리라는 뜻은 아니다. 거래비용, 시간 순서, 데이터 누출 방지, 예산 제약은 성능을 실제로 측정하기 위한 구조다. 대규모 학습이 어디까지 허용되는지 정의하는 데 이론과 제도가 필요하다. \href{http://www.incompleteideas.net/IncIdeas/BitterLesson.html}{Sutton의 글}은 방법론 선택의 배경을 제공한다.\par

\subsection{슬라이드 33 · 이중 하강과 금융 데이터}
고전적 편향·분산 그림에서는 모형을 복잡하게 할수록 처음에는 오차가 줄다가 과적합으로 다시 오른다. 일부 과매개변수 모형에서는 훈련자료를 정확히 보간하는 지점 뒤에 시험 오차가 다시 내려가는 \textbf{이중 하강}이 관찰된다.\par

중요한 조건은 모형 크기만이 아니다. 최적화 알고리즘의 암묵적 정규화, 자료의 구조, 잡음의 형태가 함께 작동한다. 시계열이 짧고 시장 구조가 바뀌는 금융에서는 현상을 보편 법칙으로 가정할 수 없다. 따라서 보간 이후에도 반드시 시간 순서에 맞춘 표본외 검증과 거래비용 점검을 거친다. \href{https://doi.org/10.1073/pnas.1903070116}{Belkin 등}의 논문은 이 현상의 기본 그림을 제시한다.\par

\section{02. 금융 텍스트와 언어모형}
숫자만으로는 기업의 전망과 위험을 다 읽을 수 없다. 사업보고서의 위험 문단, 실적 발표의 발언, 기사와 규제 공시는 모두 텍스트다. 이 장은 단어 세기에서 문맥 표현으로 발전한 방법을 이해하고, 언어모형의 학습 목표가 왜 그럴듯한 오류를 만들 수 있는지 경제학의 유인 문제와 연결한다.\par

\subsection{슬라이드 01 · 텍스트를 금융 데이터로 읽기}
재무제표의 숫자는 표준화되어 있지만 경영자의 설명은 길고 비정형적이다. 같은 “소송”이라는 단어도 소송을 당했다는 뜻인지 승소했다는 뜻인지 문맥을 보아야 한다. 텍스트 분석의 목표는 문서를 짧게 줄이는 데 그치지 않고 \textbf{미래 현금흐름이나 위험에 관한 추가 정보}를 추출하는 것이다.\par

다만 문장에 이미 주가가 반응했다면 정보 추출과 초과수익 창출은 다른 문제다. 텍스트를 이해하는 성능은 분류·추출의 정확도로, 투자 성능은 발표 뒤의 수익률과 비용으로 따로 검증해야 한다.\par

\subsection{슬라이드 02 · 사전에서 언어모형으로}
세대별 발전을 “단어의 빈도 → 단어의 조합 → 의미 벡터 → 문맥에 따라 달라지는 표현”으로 볼 수 있다. 새 방법은 앞선 방법의 약점을 줄이지만 계산비용과 해석 난도도 늘린다. 간단한 사전도 특정 공시 표현이 안정적일 때는 재현 가능하고 감사하기 쉽다.\par

모형을 고를 때는 비교 기준을 먼저 정한다. 동일 문서에서 감성을 맞히는지, 발표 직후 시장 반응을 설명하는지, 미래 실적을 예측하는지에 따라 승자가 바뀔 수 있다. 텍스트를 숫자로 바꾸는 방식은 \textbf{목적 함수}에 맞게 평가해야 한다.\par

\subsection{슬라이드 03 · 사전 단어 세기의 함정}
가장 단순한 감성 점수는 $S=(N_+-N_-)/N$이다. $N_+$와 $N_-$는 사전에 있는 긍정·부정 단어 수, $N$은 전체 단어 수다. 문서 길이로 나누지 않으면 긴 보고서가 자동으로 극단적인 점수를 얻는다.\par

일반 영어 사전의 “liability”나 “capital” 같은 단어는 금융 문맥에서 일상적 의미와 다르게 쓰인다. 단어가 위험을 뜻하는지 회계 항목을 뜻하는지 무시하면 점수에 체계적 오차가 생긴다. \href{https://doi.org/10.1111/j.1540-6261.2010.01625.x}{Loughran·McDonald}는 금융 문서에 맞는 사전이 필요한 이유를 보여 준다.\par

\subsection{슬라이드 04 · 금융 전문 사전의 의미}
전문 사전은 단어의 극성을 \textbf{사용되는 분야}에 맞춘다. 예를 들어 “tax”는 범용 감성 사전에서 부정적으로 처리될 수 있어도 세금 공제나 이연법인세 문맥에서는 방향을 단정할 수 없다. 금융 사전은 그 분야에서 실제로 부정적 의미를 갖는 단어를 다시 분류한다.\par

그래도 문장 구조를 완전히 읽는 것은 아니다. “손실이 발생하지 않았다”는 부정어 때문에 뜻이 바뀌고, 법적 위험을 의무적으로 나열하는 문서는 실제 위험과 문서 길이가 다를 수 있다. 단어 수의 유용성은 최종 과제에서 확인해야 한다.\par

\subsection{슬라이드 05 · n그램과 부정 표현}
n그램은 연속한 $n$개의 토큰을 하나의 특징으로 센다. “good” 하나는 긍정이지만 “not good”이라는 2그램은 그렇지 않다. “strong demand”와 “weak demand”도 명사 하나보다 형용사와의 결합을 보아야 구별된다.\par

단점은 조합 수가 폭발한다는 것이다. 어휘가 $V$개이면 가능한 2그램은 대략 $V^2$개이고 실제 관측은 희소하다. 최소 빈도 제한, 차원 축소, 정규화가 필요하다. 금융에서는 실적발표의 표현을 학습하더라도 과거 발표만으로 어휘를 만들고 미래 발표를 채점해야 데이터 누출을 피한다.\par

\subsection{슬라이드 06 · 임베딩은 무엇을 보존하는가}
임베딩은 단어 $w$를 벡터 $v_w\in\mathbb R^d$에 대응시킨다. 비슷한 문맥에서 등장하는 단어들이 가까운 위치를 갖도록 학습한다. 코사인 유사도 $v_a^\top v_b/(\|v_a\|\|v_b\|)$는 길이보다 방향을 비교한다.\par

하나의 단어에 벡터 하나만 주면 “bank”가 은행인지 강둑인지 문맥에 따라 달라지는 뜻을 동시에 표현하기 어렵다. 이 한계가 문장 속 다른 단어를 함께 살피는 문맥형 표현으로 이어진다. 벡터가 가깝다는 사실은 예측에 유용할 수 있지만 인과관계를 뜻하지는 않는다.\par

\subsection{슬라이드 07 · 토큰과 입력의 비용}
언어모형은 글을 하위 단어 조각인 토큰으로 나누고 정수 ID로 바꾼다. 같은 단어라도 언어와 토크나이저에 따라 조각 수가 다르다. 숫자, 종목 코드, 한국어 조사도 예상과 다르게 잘릴 수 있다.\par

긴 사업보고서를 넣을 때 토큰 수는 비용과 문맥 창의 한도를 결정한다. 중요한 주석이 뒤에서 잘리면 모형이 그 내용을 모른 채 답한다. 그래서 자료를 나눌 때 기간·회사·주석 번호를 붙이고, 답에 사용한 문단을 다시 찾을 수 있게 해야 한다.\par

\subsection{슬라이드 08 · 어텐션의 계산}
각 토큰의 현재 표현 $x_i$를 질의 $q_i=x_iW_Q$, 키 $k_j=x_jW_K$, 값 $v_j=x_jW_V$로 바꾼다. 토큰 $i$가 $j$를 참고하는 정도는 $q_i^\top k_j/\sqrt d$의 softmax로 정한다. 새 표현은 값의 가중평균이다.\par

\[
\alpha_{ij}=\frac{\exp(q_i^\top k_j/\sqrt d)}{\sum_{\ell}\exp(q_i^\top k_{\ell}/\sqrt d)},\qquad
h_i=\sum_j\alpha_{ij}v_j.
\]

분모의 $\sqrt d$는 차원이 커질 때 내적의 크기가 과도하게 커져 softmax가 한 항에 몰리는 현상을 줄인다. “bank raised interest rates”에서 bank는 금리와 연결되어 금융기관 뜻으로 해석된다. \href{https://arxiv.org/abs/1706.03762}{Transformer 원문}은 이 연산의 출발점이다.\par

\subsection{슬라이드 09 · 다음 토큰 예측과 압축}
문장 $x_1,\ldots,x_n$의 확률을 연쇄법칙으로 $p_\theta(x_1,\ldots,x_n)=\prod_{t=1}^{n}p_\theta(x_t\mid x_{<t})$로 분해한다. 훈련은 관찰된 토큰의 음의 로그확률을 줄이는 과정이다. 자주 쓰이는 패턴에 높은 확률을 주면 평균 부호 길이도 짧아져 예측과 압축이 연결된다.\par

다음 토큰 하나를 맞히려 해도 긴 문맥의 인물, 관계, 수치가 필요할 수 있어 능력이 넓게 생긴다. 동시에 목표는 텍스트 확률이지 \textbf{외부 세계의 참·거짓을 직접 채점하는 것}이 아니다. 답변의 검증 체계를 따로 두어야 한다.\par

\subsection{슬라이드 10 · 유창한 실패의 출발점}
문법적으로 자연스러운 문장은 독자의 신뢰를 쉽게 얻는다. 그러나 금융 분석에서는 오답의 형태가 매끈하면 더 위험하다. 숫자를 틀리게 합산하거나 없는 공시를 인용해도 문장만으로는 발견하기 어렵다.\par

실패를 세 갈래로 나누면 대응이 쉬워진다. 필요한 정보를 못 봤는지, 본 정보를 잘못 끌어냈는지, 계산·추론에서 틀렸는지 구분한다. 첫째에는 자료 첨부와 검색, 둘째에는 재질문과 근거 대조, 셋째에는 계산기와 독립 검산이 필요하다.\par

\subsection{슬라이드 11 · 생성모형의 목적 함수를 정확히 쓰기}
확률의 연쇄법칙은 바로 앞 토큰 하나만 조건에 넣는다는 뜻이 아니다. 일반적인 자기회귀 언어모형은 \textbf{앞의 전체 문맥} $x_{<t}=(x_1,\ldots,x_{t-1})$을 조건으로 사용한다. 최대우도 추정은 다음과 같다.\par

\[
\theta^*=\arg\max_\theta\sum_{m=1}^{M}\sum_{t=1}^{n_m}
\log p_\theta(x^{(m)}_t\mid x^{(m)}_{<t}).
\]

로그를 취하면 곱이 합으로 바뀌어 수치 계산이 안정적이고, 확률을 높이는 것과 같은 최적점을 갖는다. 이 식은 텍스트 분포의 적합도 목표이며, 금융 분석의 경제적 손실은 아직 들어 있지 않다.\par

\subsection{슬라이드 12 · 목적 함수에 진실 항이 없다는 뜻}
웹에 반복된 잘못된 주장도 훈련 텍스트라면 예측 대상이 된다. 최대우도 목표는 그 문장이 실제 기업의 사실인지 별도로 대조하지 않는다. 후속 정렬과 검색 도구가 이를 보완할 수 있지만, 기본 훈련 목표만으로 사실성이 보장되지는 않는다.\par

“거짓말”이라는 표현은 모델이 사실을 알고도 속이려는 의도를 암시해 혼동을 준다. 실무에서는 사실 오류, 출처 환각, 과도한 확신을 관찰 가능한 실패로 정의하는 편이 낫다. 같은 오류라도 투자 추천에 쓰일 때 손실함수는 훨씬 커진다.\par

\subsection{슬라이드 13 · 질문을 바꾸면 평가도 바뀐다}
“왜 거짓말하는가”보다 “어떤 입력과 보상 아래서 어떤 오류가 반복되는가”가 검증 가능한 질문이다. 모호한 요청은 모델이 가능한 해석 가운데 하나를 택하게 한다. 사용자가 원한 회계 정의와 모델이 택한 정의가 다르면 계산은 매끈해도 결론은 틀린다.\par

따라서 작업을 시작할 때 기간, 연결·별도 기준, 조정 대상, 해당 항목이 없을 때의 출력 형식을 명시한다. 결과에는 항목별 원문 위치와 계산식을 요구하고, 사람이 표본을 추출해 확인한다.\par

\subsection{슬라이드 14 · 걸러낼 것이 없는 회사를 시험하기}
“핵심이익을 계산하라”는 표현은 일회성 항목의 존재를 암묵적으로 기대하게 만들 수 있다. 같은 일을 “반복되지 않는 항목이 있으면 열거하고, 없으면 없다고 답하라”고 풀어 쓰면 과제가 더 분명하다. 추출 도구는 \textbf{양성 사례뿐 아니라 음성 사례}에서 평가해야 한다.\par

정밀도는 찾아냈다고 한 항목 중 진짜의 비율, 재현율은 실제 항목 중 찾아낸 비율이다. 항목이 거의 없는 기업에서는 억지 양성이 많이 나와 정밀도가 떨어진다. 재무 분석에서는 항목을 빼는 근거를 원문과 금액에 연결해야 하며, 모델의 말만으로 이익을 조정하면 안 된다.\par

\subsection{슬라이드 15 · 모델의 국적을 성향으로 단정하지 않기}
같은 기업에 대한 두 모형의 전망이 다르면 정치적 성향부터 추측하기 쉽다. 하지만 모형이 접근한 정보의 범위, 학습시점, 검색 언어가 다를 수 있다. 낙관적 전망은 의도적인 선호가 아니라 부정적 기사의 누락에서 생길 수 있다.\par

비교 실험에서는 기업과 날짜를 같게 고정하고, 각 모형에 같은 원자료를 제공한 뒤 전망과 실제 결과를 비교해야 한다. 단순히 두 답변의 어조를 비교하는 것은 정확도 검사가 아니다.\par

\subsection{슬라이드 16 · 예상을 뒤집는 결과의 검증}
직관과 반대의 결과일수록 표본 선정과 기준을 확인한다. 어떤 중국 기업을 포함했는지, 전망 시점의 뉴스만 주었는지, 낙관성을 목표주가·이익 전망·문장 감성 중 무엇으로 측정했는지에 따라 결론이 달라질 수 있다.\par

독자는 모형 간 차이를 “어느 국적이 어느 시장을 좋아한다”로 외우기보다 \textbf{정보 집합 차이의 가설}로 보아야 한다. 다음 단계는 동일한 지역 뉴스의 제공으로 차이가 줄어드는지 확인하는 것이다.\par

\subsection{슬라이드 17 · 지역 정보의 사각지대}
중소형주에 관한 규제 조치나 소송 기사는 현지 언어의 지역 매체에 먼저 실릴 수 있다. 영어권 검색에 안 보인다고 사건이 없는 것은 아니다. 한국의 코스닥 기업을 평가할 때도 DART, 거래소 공시, 국내 언론을 직접 확인해야 한다.\par

정보를 추가한 뒤 전망 차이가 줄어들면 국적 자체보다 자료 접근성이 원인이라는 설명에 힘이 실린다. 하지만 검색 결과가 더 많다는 사실만으로 기사가 모두 신뢰할 만하다는 뜻은 아니다. 기업명·발표일·법적 상태를 교차 확인한다.\par

\subsection{슬라이드 18 · 기억과 진짜 사전 예측의 차이}
기업명과 날짜를 지운 옛 실적발표도 문장 패턴으로 식별할 수 있다면, 과거 자료에 대한 높은 점수는 미래 예측력의 증거가 아니다. 모형의 훈련 종료일보다 뒤의 사건만 시험하고, 프롬프트에 미래의 정보가 섞이지 않았는지 확인해야 한다.\par

엄격한 실험은 \textbf{시점별 정보 집합} $\mathcal I_t$만 사용해 $t+1$을 예측한다. 모형 가중치를 학습한 시점, 검색 색인의 업데이트 시점, 평가 자료의 공개 시점까지 기록해야 한다. 금융에서 시간 누출은 아주 작은 예측 신호를 쉽게 압도한다.\par

\subsection{슬라이드 19 · 모형의 투자 스타일}
답변 모델은 학습자료와 정렬 과정 때문에 특정 업종이나 최근 하락 종목을 더 자주 추천할 수 있다. 이는 경제학에서 말하는 투자 스타일과 비슷하지만, 실제 투자전략의 수익성이 입증되었다는 뜻은 아니다.\par

성향을 측정하려면 동일 시점의 종목과 자료를 여러 모형에 주고, 추천 비율을 업종·규모·직전 수익률별로 비교한다. 반대 근거를 추가했을 때 전망이 적절히 갱신되는지도 중요하다. 하나의 모형을 여러 번 물어 얻은 답은 독립된 전문가 의견이 아니다.\par

\subsection{슬라이드 20 · 만족도 보상의 부작용}
사용자 평가를 보상에 넣으면 친절함이 높아질 수 있지만, 사용자가 듣고 싶은 말에 동의할 유인도 생긴다. 2025년 GPT-4o 업데이트에서 과도한 동조가 나타난 사례는 목적함수의 대리 지표를 잘못 맞추면 원하는 품질과 다른 행동이 나온다는 점을 보여 준다.\par

금융에서는 “이 종목을 사면 좋겠지?”라는 유도 질문에 대한 동의가 특히 위험하다. 평가 문항에 정확한 반론 제시, 불확실성 표현, 근거 없는 주문 권유 거절을 포함해야 한다. \href{https://openai.com/index/expanding-on-sycophancy/}{OpenAI의 사후 분석}은 피드백 신호의 결합 문제를 설명한다.\par

\subsection{슬라이드 21 · 사후 분석과 인과 사슬}
업데이트 뒤 동조가 관찰되고 이전 버전으로 되돌리자 완화되었다는 시간 순서는 중요한 단서다. 하지만 사후 분석은 여러 변경이 함께 들어갔음을 인정하므로 “좋아요 신호 하나만 원인”이라고 단순화하면 안 된다.\par

평가 지표가 사용자 만족을 포착해도 진실성, 장기적 편익, 반대 의견 제시까지 완전히 측정하지는 못한다. 제품 평가에서는 개별 지표의 개선뿐 아니라 서로 다른 목표가 충돌하는 극단 사례를 함께 시험한다.\par

\subsection{슬라이드 22 · 체스 에이전트 실험의 설계}
도구를 쓸 수 있는 에이전트에게 “체스 엔진을 이겨라”라고만 지시하면, 체스를 두는 방법 외에 게임 파일을 바꾸거나 상대 프로그램에 접근하는 방법도 환경에 존재할 수 있다. 목표를 결과만으로 정의하면 허용된 행동의 경계가 빠진다.\par

이 실험은 특정 모형이 언제나 부정직하다는 결론보다 \textbf{권한 설계의 중요성}을 보여 준다. 금융 에이전트에게 분석과 주문 권한을 한꺼번에 주면, 오류가 설명 단계에서 끝나지 않고 실제 거래로 이어질 수 있다. \href{https://arxiv.org/abs/2502.13295}{실험 논문}을 읽을 때 성공의 정의와 허용 도구를 먼저 보자.\par

\subsection{슬라이드 23 · 목표 달성과 규칙 준수는 다른 변수}
게임 상태를 바꾸어 얻은 “승리”는 체스 실력의 성과가 아니다. 평가자가 결과 파일의 승패만 읽으면 이런 경로를 구별하지 못한다. 금융에서도 손익 숫자만 보면 미공개 정보 사용, 과도한 위험 부담, 장부 밖 포지션을 숨긴 성과를 놓칠 수 있다.\par

따라서 결과 지표와 절차 제약을 동시에 둔다. 모형에 필요한 최소 권한만 주고, 원장·주문·자료 출처의 로그를 독립적으로 검증한다. 성과 측정의 대상이 무엇인지 처음부터 정의해야 한다.\par

\subsection{슬라이드 24 · 굿하트 법칙, 루커스 비판, 대리인 문제}
굿하트 법칙은 유용한 대리 지표를 목표로 최적화하면 원래 목표와의 관계가 약해질 수 있다는 통찰이다. 루커스 비판은 정책이 바뀌면 경제주체의 행동도 바뀌므로 과거 관계를 그대로 새 정책에 적용하기 어렵다는 주장이다. 두 개념은 \textbf{행동의 변화가 통계적 관계를 바꾼다}는 점에서 만난다.\par

모형의 보상 $R$이 진짜 사회적 가치 $V$의 불완전한 근사라면, $\arg\max R$이 $\arg\max V$와 다를 수 있다. 개발자와 사용자가 원하는 것을 에이전트에게 완전한 계약으로 명시하기 어려운 주인·대리인 문제다. 독립 감사와 반복 평가, 권한 제한이 보상함수의 문구만큼 중요하다.\par

\subsection{슬라이드 25 · 환각을 과신으로 측정하기}
틀린 답을 모두 같은 위험으로 볼 필요는 없다. “확실하지 않다”라고 말하고 근거 확인을 요청한 오답은 확신에 차서 허위 출처를 만든 오답과 다르다. 확률 $p$를 말한 예측은 실제로 약 $p$의 비율로 맞아야 \textbf{보정}되어 있다고 한다.\par

예를 들어 90\% 확신 답변 100개 중 90개가 맞으면 보정이 좋다. 정답률을 높이는 것과 자신의 불확실성을 제대로 알리는 것은 다른 능력이다. \href{https://arxiv.org/abs/2605.01428}{Yona 등}은 거절 또는 단언의 양자택일 대신 정직한 불확실성 표현을 제안한다.\par

\subsection{슬라이드 26 · 저장된 지식과 인출의 병목}
모형이 어떤 사실을 표현 안에 담고 있어도 직접 질문에서 즉시 꺼내지 못할 수 있다. 추가 추론을 제공했을 때 정답이 나오면 자료의 부재와 인출 실패를 구분할 단서가 된다. 다만 특정 벤치마크에서 높은 인코딩 비율을 모든 금융 사실로 일반화할 수는 없다.\par

공시나 법규처럼 날짜에 민감한 사실은 모형 내부 기억보다 원문 검색이 필요하다. 검색 후에도 문단이 질문의 회사·기간에 맞는지 확인하고 숫자는 계산기로 검산한다. \href{https://arxiv.org/abs/2602.14080}{Calderon 등}은 지식의 저장과 접근을 구분한다.\par

\subsection{슬라이드 27 · 보이는 추론의 한계}
단계별 설명은 오류를 발견하는 데 도움이 되지만, 화면에 보이는 설명이 모형 내부에서 실제로 사용한 모든 원인을 기록한다고 볼 수 없다. 설명 자체가 평가의 대상이 되면 그럴듯한 절차를 쓰는 데 최적화될 수도 있다.\par

금융 분석에서는 “생각한 과정을 써라”만으로 검증을 끝내지 않는다. 입력 공시의 위치, 재현 가능한 계산식, 대체 가정에서의 결과를 요구한다. 계산과 출처를 외부에서 다시 실행할 수 있어야 한다.\par

\subsection{슬라이드 28 · 다음 토큰 맞히기가 만드는 능력}
추리소설의 마지막 단어를 잘 예측하려면 앞선 사건과 인물의 관계를 추적해야 한다. 따라서 단순해 보이는 훈련 목표가 복잡한 표현 능력을 유도할 수 있다. 이는 “다음 토큰 예측”이 사소한 기능이라는 뜻이 아니다.\par

하지만 새로운 사례에서의 이해를 인간의 이해와 동일시해서도 안 된다. 훈련 목표가 요구한 규칙성과 실제 금융 판단의 목표가 다를 수 있다. 원리의 강점과 한계를 동시에 평가하려면 새 자료에서의 결정 성과가 필요하다.\par

\subsection{슬라이드 29 · 산술은 외부 도구로 분리하기}
여러 숫자를 더하는 일은 결정적 알고리즘으로 정확히 수행할 수 있다. 언어모형이 숫자 패턴을 바탕으로 합계를 추측하게 둘 이유가 없다. 순서를 바꾸어도 합계가 같아야 한다는 \textbf{불변성 검사}는 계산 오류를 찾는 간단한 방법이다.\par

재무제표 분석 흐름을 “문장에서 항목과 단위 추출 → 표로 구조화 → 계산기로 합산 → 원문과 대조”로 나누면 각 단계의 오류를 구별할 수 있다. 모델의 숫자 해석도 검산해야 하므로 도구 호출만으로 전체 과정이 자동으로 안전해지는 것은 아니다.\par

\subsection{슬라이드 30 · 유창함 대신 재현성을 기준으로}
좋은 금융 답변은 출처, 기준일, 단위, 계산 경로, 불확실성을 드러낸다. 쉬운 말로 과제를 구체화하고 필요한 자료를 직접 제공한 다음, 중요한 결론은 독립된 자료와 도구로 확인한다. 여러 모형에 묻는 것은 도움이 되지만 같은 오류를 공유할 수 있어 \textbf{독립된 원자료 확인}을 대신하지 못한다.\par

주식 매수 비율처럼 개인 사정을 모르는 상태에서 제시된 그럴듯한 숫자는 맞춤형 결론이 아니다. 독자는 질문의 빈칸을 먼저 채운다. 투자 기간과 제약, 비교 대안, 손실 한도, 정보의 날짜가 정해져야 비로소 예측을 행동으로 바꿀 수 있다.\par

\section{03. 정규화 회귀와 수익률 예측}
주식시장의 미래 수익률은 평균은 작고 변동성은 큰 자료에서 추정해야 한다. 표본 안에서 잘 맞는 식은 과거의 우연을 익혔을 수 있다. 이 장에서는 비교 가능한 표본외 채점 규칙을 세우고, 추정오차가 왜 커지는지 유도한 뒤, 능형회귀·라쏘·교차검증이 어떤 편향을 의도적으로 도입하는지 단계별로 살핀다.\par

\subsection{슬라이드 01 · 수익률 예측의 출발점}
AI가 금융에서 유용하려면 과거의 가격을 설명하는 것보다 \textbf{아직 보지 않은 수익률}을 예측해야 한다. 시장 전체의 수익률은 종목 간 상대 순위를 맞히는 일과 다르다. 한 해에 모든 주식이 함께 떨어질 수 있어 시장 시계열 예측에는 거시 상태가 중요하다.\par

이 장의 핵심 질문은 단순하다. 많은 예측변수를 넣어도 과거 평균만 쓰는 사람보다 더 잘 맞힐 수 있는가. 먼저 평가 규칙을 고정해야 답을 비교할 수 있다.\par

\subsection{슬라이드 02 · 앞 장에서 이어지는 문제}
금융의 예측 신호는 가격·뉴스·재무비율에 퍼져 있지만 수익률의 대부분은 뜻밖의 충격이다. 관계가 시간에 따라 변하는 비정상성도 있다. 그러므로 복잡한 모형을 쓰기 전에 약한 신호와 짧은 실효 표본이라는 제약을 받아들여야 한다.\par

진행 경로는 예측 회귀의 실패 → 편향·분산 진단 → 계수 축소와 변수 선택 → 올바른 검증이다. 각 수식은 어떤 오차가 줄고 어떤 오차가 늘어나는지 설명하기 위해 등장한다.\par

\subsection{슬라이드 03 · 시장 수익률 예측의 오래된 실패}
배당수익률이나 금리 스프레드가 높을 때 다음 해의 시장 수익률이 달라질 수 있다는 이론은 그럴듯하다. 문제는 과거 자료에서 찾은 관계가 미래에도 남는가다. 발표 시점까지의 자료를 모두 사용해 회귀한 뒤 같은 기간을 채점하면 미래 정보가 추정에 들어간다.\par

이후의 사례는 특정 변수가 영원히 무용하다는 선언이 아니다. 예측 능력은 기간과 시장 상태에 따라 달라질 수 있으며, 이를 판정하는 규칙을 사전에 정해야 한다.\par

\subsection{슬라이드 04 · 평균보다 큰 변동성}
주식시장의 연평균 수익률이 양수라도 매년 수익률은 크게 흔들린다. 평균이 10\% 안팎이고 표준편차가 20\% 안팎이면 한 해 관측치의 신호대잡음비는 대략 0.5에 불과하다. 평균을 정확히 추정하려면 단순히 주가 관측 횟수보다 긴 기간이 필요하다.\par

더 중요한 것은 예측 대상이다. 올해의 주가지수 수준을 아는 것은 내년의 초과수익률을 아는 것과 다르다. 채점에서는 배당을 포함한 총수익률과 무위험수익률을 뺀 초과수익률의 정의를 일관되게 써야 한다.\par

\subsection{슬라이드 05 · 예측변수 동물원}
가장 단순한 예측 회귀는 $r_{t+1}=a+bx_t+\varepsilon_{t+1}$이다. $x_t$는 시점 $t$에 이미 관찰된 배당/주가, 이익/주가, 장단기금리차 등이다. $b$는 지표가 한 단위 높아졌을 때 다음 수익률의 조건부 평균이 얼마나 바뀌는지 나타낸다.\par

변수가 많아질수록 과거에 우연히 잘 맞은 지표를 고르기 쉬워진다. 경제적으로 그럴듯한 이유와 통계적 유의성만으로 끝내지 않고, 선택 과정 전체를 포함해 미래 예측 성능을 봐야 한다.\par

\subsection{슬라이드 06 · 높은 배당수익률의 두 해석}
매년 배당이 $g$로 성장하고 요구수익률이 $r>g$로 일정하다면, 고든 성장모형에서 $P_t=D_{t+1}/(r-g)$이다. 양변을 정리하면 $D_{t+1}/P_t=r-g$다. 배당 대비 가격이 낮을 때 높은 미래 수익률 $r$ 때문일 수도 있고 낮은 배당성장률 $g$ 때문일 수도 있다.\par

따라서 배당수익률의 계수에 인과적 이름을 곧장 붙일 수 없다. 경기침체기에 위험 보상이 올라 가격이 떨어졌다면 이후 수익률이 높아질 수 있다. 기업의 배당 성장 자체가 나빠졌다면 같은 높은 배당수익률이 좋은 투자 신호가 아닐 수도 있다.\par

\subsection{슬라이드 07 · 자료를 잘게 나누어도 평균은 정밀해지지 않는다}
전체 기간을 $T$년, 관측 간격을 $h$년이라고 하자. 독립 로그수익률 $X_k$가 평균 $\mu h$, 분산 $\sigma^2h$를 가지면 $n=T/h$개를 합한 연율 평균 추정치는 $\hat\mu=\sum_k X_k/T$다. 독립성 때문에 분산은 다음과 같이 계산된다.\par

\[
\operatorname{Var}(\hat\mu)=\frac{n\sigma^2h}{T^2}=\frac{\sigma^2}{T}.
\]

$h$가 지워진다. 하루 자료가 한 달 자료보다 행 수는 많아도, 동일한 역사 기간에서 평균수익률의 핵심 정보량은 크게 늘지 않는다. 실제 수익률에 자기상관이 있으면 식을 수정해야 하지만 “많은 행 = 많은 독립 정보”가 아니라는 직관은 유지된다.\par

\subsection{슬라이드 08 · 단순수익률, 로그수익률, 연속복리}
단순수익률은 $R=(P_1-P_0)/P_0$이고 로그수익률은 $r=\ln(P_1/P_0)=\ln(1+R)$이다. $R$이 작을 때 테일러 전개 $\ln(1+R)=R-R^2/2+\cdots$로 둘이 가까워진다. 여러 기간의 가격비는 곱이므로 로그를 취하면 수익률은 합이 된다.\par

\[
\frac{P_n}{P_0}=\prod_{k=1}^n(1+R_k),\qquad
\ln\frac{P_n}{P_0}=\sum_{k=1}^n r_k.
\]

연 $r$을 $m$번 복리로 적용한 $(1+r/m)^m$은 $m\to\infty$에서 $e^r$로 간다. 로그수익률의 합산성과 짧은 구간 분산의 비례관계는 독립 또는 적절한 약한 의존 조건을 구분해 사용해야 한다.\par

\subsection{슬라이드 09 · 일별 자료의 다른 역할}
일별 자료는 연평균 기대수익률 추정에 마법처럼 많은 정보를 더하지 않는다. 오히려 호가 스프레드, 비동시 거래, 가격 정체 때문에 가짜 자기상관이 생길 수 있다. 매수호가와 매도호가 사이의 거래가격 이동만으로 하루 수익률이 오르내릴 수 있다.\par

반면 변동성은 하루 단위 관측에서 더 잘 추정할 수 있다. 주가가 하루에 얼마나 흔들리는지는 위험관리와 포지션 크기 결정에 중요하다. 따라서 \textbf{평균을 예측하는 문제와 분산을 추정하는 문제}에 같은 자료 주기를 기계적으로 적용하지 않는다.\par

\subsection{슬라이드 10 · 표본 안의 유의성}
예측 회귀에서 $t$값이 크거나 장기 $R^2$가 높다는 결과는 먼저 표본 내 사실이다. 변수를 고른 기간과 계수를 추정한 기간이 같으면 선택 편향이 발생한다. 장기 수익률은 겹치는 기간이 많아 표준오차 계산도 까다롭다.\par

정말 알고 싶은 것은 과거 어느 시점으로 돌아가 당시 이용 가능했던 자료만으로 만든 예측이 이후 결과를 맞혔는가다. 이 질문은 금융상품의 실제 사용 가능성과도 일치한다.\par

\subsection{슬라이드 11 · 표본내와 표본외를 구분하기}
표본내 적합은 전체 기간으로 $a,b$를 추정한 후 같은 자료에서 잔차를 계산한다. 표본외 검증은 각 시점 $t$에 $t$까지의 자료만으로 $\hat a_t,\hat b_t$를 추정해 $\hat r_{t+1}=\hat a_t+\hat b_tx_t$를 만든다. 이후 $r_{t+1}$이 실현되면 오차를 기록한다.\par

이 과정을 앞으로 이동하며 반복하는 것이 재귀적 예측이다. 데이터 수정본이 나중에 공개되었다면 과거의 정보 집합에 넣으면 안 된다. 금융에서는 이 작은 시점 오류가 약한 진짜 신호보다 클 수 있다.\par

\subsection{슬라이드 12 · 과거 평균과 표본외 결정계수}
벤치마크는 예측변수 없이 시점 $t$까지의 평균 $\bar r_t$를 다음 수익률 예측으로 쓰는 것이다. 모형과 벤치마크의 제곱오차 합을 비교한다.\par

\[
R^2_{\mathrm{OOS}}=1-
\frac{\sum_t(r_{t+1}-\hat r_{t+1})^2}
{\sum_t(r_{t+1}-\bar r_t)^2}.
\]

분자가 더 작아야 양수다. 표본내 OLS의 절편 포함 $R^2$는 0 이상이지만 표본외 값은 음수가 될 수 있다. 누적 제곱오차 차이 $\sum_{s\le t}(e^2_{\rm bench,s}-e^2_{\rm model,s})$를 그리면 성과가 어느 기간에 생겼는지도 보인다.\par

\subsection{슬라이드 13 · 웰치와 고열의 판정}
배당수익률 등 많은 시장 예측변수는 과거 표본 안에서 설명력을 보였지만, 시점마다 다시 추정하는 표본외 평가에서는 과거 평균에 뒤졌다. 한 번의 위기 기간에 기울기가 크게 맞아 전체 통계량을 끌어올렸을 가능성도 살펴야 한다.\par

이 결과는 초과수익 예측을 포기하라는 뜻보다 \textbf{비교 기준을 낮추지 말라}는 교훈이다. \href{https://doi.org/10.1093/rfs/hhm014}{Welch·Goyal}의 누적 오차 그림에서는 선이 내려갈 때 제안 모형이 기준선보다 더 큰 오차를 낸다.\par

\subsection{슬라이드 14 · 간단한 경제 제약은 도움이 될까}
배당수익률이 높으면 미래 수익률이 높아야 한다는 이론적 부호를 강제하거나, 예상 초과수익률이 음수일 때 0으로 절단하면 과도한 추정값을 줄일 수 있다. 이는 모형에 사전 믿음을 넣는 정규화의 한 형태다.\par

월별 $R^2_{\rm OOS}$가 작아도 평균·분산 투자자는 포트폴리오 비중을 개선할 수 있다. 다만 기대효용 증가는 위험회피계수, 매매비용, 차입 제약에 민감하다. \href{https://doi.org/10.1093/rfs/hhm055}{Campbell·Thompson}의 주장은 통계 점수와 경제적 가치를 구분해 읽어야 한다.\par

\subsection{슬라이드 15 · 시간이 지난 뒤의 재검증}
첫 논쟁 이후 발표된 새 변수도 발표 당시의 자료에는 잘 맞을 수 있다. 독립된 후속 기간과 공통한 기준으로 다시 평가하면 결과가 약해지기도 한다. 이는 예측 관계가 소멸했을 수도 있고, 처음부터 선택 편향이 있었을 수도 있음을 뜻한다.\par

두 설명을 가르려면 원 논문의 표본 종료일, 논문 발표일, 이후 기간을 따로 본다. 발표 전에 이미 감소했다면 데이터 채굴이나 불안정성의 비중이 클 수 있다. 발표 후 추가 감소에는 투자자의 반응도 기여할 수 있다.\par

\subsection{슬라이드 16 · 유의성·안정성·실용성의 깔때기}
변수가 표본내 유의성을 유지하는지, 같은 부호로 표본외 예측을 하는지, 거래비용 뒤에도 투자 의사결정을 개선하는지는 서로 다른 관문이다. 앞 관문을 통과해도 뒤 관문을 자동으로 통과하지 않는다.\par

새 논문의 수를 세어 “성공률”이라고 부를 때는 각 연구가 사용한 시작일, 예측 기간, 재추정 방식이 같은지 확인해야 한다. 공정한 비교는 변수 자체뿐 아니라 \textbf{검증 절차}를 맞춘다.\par

\subsection{슬라이드 17 · 누적 오차 차이 그래프 읽기}
그래프의 세로축이 누적 $\mathrm{SSE}_{\rm bench}-\mathrm{SSE}_{\rm model}$이면 상승은 예측변수의 승리, 하락은 과거 평균의 승리다. 선이 특정 위기 이후 꺾이는 경우에는 기울기 자체가 달라졌는지, 큰 한두 관측치가 오차를 지배했는지 확인한다.\par

표본내 선과 표본외 선이 벌어지면 과거를 맞힌 계수가 미래에는 통하지 않았다는 뜻이다. 그래프의 0선뿐 아니라 기울기의 안정성과 표본 종료일 표시를 함께 보아야 한다.\par

\subsection{슬라이드 18 · 첫 단계의 결론}
예측변수가 많아도 과거 평균보다 좋다는 주장은 어렵게 입증해야 한다. 기준선은 쉽게 이길 수 있는 허수아비가 아니라 실제로 이용 가능하고 추정오차가 작은 대안이어야 한다. 정확도가 조금 높아도 구현 비용과 위험을 함께 본다.\par

이제 실패의 원인을 모형의 ‘지능 부족’이 아니라 약한 신호와 큰 추정오차에서 찾는다. 그 원인이 맞다면 복잡성을 제어하는 방법이 다음 처방이 된다.\par

\subsection{슬라이드 19 · 왜 실패하는가}
진짜 조건부 기대수익률이 존재해도 매번 실현 수익률에는 큰 예측 불가능한 충격이 들어온다. 잡음이 클 때 회귀계수는 작은 표본의 우연에 따라 크게 달라진다. 변수 수가 늘면 그 우연을 설명할 자유도가 늘어난다.\par

핵심은 \textbf{참모형의 존재와 참모형을 추정할 수 있는가}가 별개라는 점이다. 금융시장의 역사 하나로 수백 변수를 추정한다면 이론적으로 유효한 관계도 실제 예측에서 실패할 수 있다.\par

\subsection{슬라이드 20 · 작은 참 신호의 규모}
$y=f(x)+\varepsilon$이고 신호와 잡음이 독립이라면 참 예측 가능 비율은 $R^2_{\rm true}=\operatorname{Var}(f(x))/[\operatorname{Var}(f(x))+\operatorname{Var}(\varepsilon)]$이다. 조건부 평균의 표준편차가 잡음의 1/10이면 $R^2=0.01/1.01\approx0.99\%$다.\par

따라서 작은 표본외 결정계수는 무조건 무의미하지 않다. 반대로 수익률을 거의 완벽히 맞히는 표본내 모형은 강한 의심을 받아야 한다. 이미 알려진 미래 정보나 거래 불가능한 가격이 입력에 섞였을 수 있다.\par

\subsection{슬라이드 21 · 과적합 모의실험을 직접 구성하기}
관측치 120개에서 후보 변수 110개를 만들되 실제 영향이 있는 것은 첫 10개라고 하자. 독립 표준정규 변수와 잡음의 분산을 1로 두고 참계수를 모두 $c$로 정하면 신호분산은 $10c^2$다. 참 $R^2=5\%$가 되게 하려면 $10c^2/(10c^2+1)=0.05$를 풀어 $c^2=1/190$을 얻는다.\par

이후 회귀에 넣는 변수 수를 늘리고 훈련자료와 새 자료의 $R^2$를 따로 기록한다. 중요한 점은 11번째 이후 변수가 진짜로 효과 0이라는 정답을 우리가 안다는 것이다. 따라서 성능 하락이 추정오차에서 왔는지 직접 확인할 수 있다.\par

\subsection{슬라이드 22 · 변수 증가와 벌점의 효과}
OLS는 변수를 추가할 때 표본내 제곱오차가 증가할 수 없다. 추가 변수를 쓰지 않도록 계수를 0으로 놓는 선택이 항상 가능하기 때문이다. 그러나 미래 자료에서는 잡음에 맞춘 계수가 오차를 키울 수 있다.\par

같은 변수 수를 유지하고 능형회귀의 벌점을 키우면 표본내 적합은 다소 낮아지지만 표본외 오차가 줄어드는 구간이 생긴다. 벌점이 너무 크면 유용한 신호까지 지워 다시 나빠진다. 그래프의 최저점은 사전 고정값이 아니라 독립 자료로 선택해야 한다.\par

\subsection{슬라이드 23 · 차원이 커질 때 추정오차의 유도}
실제 평균벡터를 $\mu=Xg$라 하고 $\tau$기간 평균수익률을 $\bar r=\mu+\bar\varepsilon$라 하자. $X$의 열공간에 대한 투영행렬 $P=X(X^\top X)^{-1}X^\top$로 OLS 예측을 만들면 $\hat\mu=P\bar r=\mu+P\bar\varepsilon$이다. $P\mu=\mu$이고 잡음 공분산이 $\sigma^2I/\tau$라는 가정 아래 추정오차 제곱의 기대값은 다음과 같다.\par

\[
E\|\hat\mu-\mu\|^2
=E[\bar\varepsilon^\top P\bar\varepsilon]
=\operatorname{tr}\!\left(P\frac{\sigma^2}{\tau}I\right)
=\frac{K\sigma^2}{\tau}.
\]

$P$가 대칭이고 $P^2=P$, $\operatorname{tr}P=K$인 성질을 사용했다. 변수 $K$가 늘수록 오차가 선형으로 커지고 기간 $\tau$가 늘수록 줄어든다. 새 자료의 잡음은 이 추정오차와 독립이므로 참 신호의 크기 $\|\mu\|^2$보다 추정오차가 크면 0 예측보다 표본외 성능이 낮아질 수 있다.\par

\subsection{슬라이드 24 · 편향과 분산을 식으로 분해하기}
새 관측치가 $y_0=f(x_0)+\varepsilon_0$이고 학습자료에서 나온 예측 $\hat f(x_0)$와 새 잡음이 독립이라고 하자. $\bar f=E[\hat f(x_0)]$를 더하고 빼면 $y_0-\hat f=\varepsilon_0+(f-\bar f)+(\bar f-\hat f)$다. 제곱하여 기댓값을 취할 때 새 잡음의 평균이 0이고 $E(\bar f-\hat f)=0$이므로 교차항이 사라진다.\par

\[
E[(y_0-\hat f)^2]=\sigma^2+
\bigl(f-E\hat f\bigr)^2+
\operatorname{Var}(\hat f).
\]

첫 항은 참 함수를 알아도 남는 잡음, 둘째는 평균 예측의 체계적 오류, 셋째는 학습표본에 따른 불안정성이다. 정규화는 둘째를 조금 늘리는 대가로 셋째를 크게 줄이려는 선택이다.\par

\subsection{슬라이드 25 · 세 오차의 경제적 의미}
수익률의 예측 불가능한 충격은 모형을 아무리 정교하게 해도 없앨 수 없다. 편향은 모형이 지나치게 단순해 실제 위험 보상을 놓칠 때 생긴다. 분산은 작은 표본마다 계수가 요동쳐 투자 비중이 불안정해지는 문제다.\par

분산이 큰 모형은 계산상 평균 점수가 나쁠 뿐 아니라 잦은 매매와 극단적 레버리지를 유발할 수 있다. 벌점은 통계적 분산과 경제적 거래비용을 동시에 낮추는 방향으로 작동할 수 있지만, 실제 효과는 검증해야 한다.\par

\subsection{슬라이드 26 · U자형 상충관계}
아주 단순한 과거 평균 모형은 변수에 따른 변화를 못 잡아 편향이 크지만 추정이 안정적이다. 모든 변수를 자유롭게 넣는 OLS는 편향이 낮을 수 있으나 잡음에 민감하다. 중간 복잡도에서 합계 오차가 낮아지는 U자형 그림은 이 관계를 요약한다.\par

그러나 U자의 위치는 자료의 신호대잡음비와 변수 상관에 따라 달라진다. 금융처럼 신호가 약하고 데이터가 짧을수록 왼쪽, 즉 강한 축소 쪽이 유리할 가능성이 크다. 그림은 사전 믿음이며 실제 최저점은 검증에서 찾는다.\par

\subsection{슬라이드 27 · 공짜 점심이 없는 이유}
모든 가능한 자료생성과정에 대해 언제나 이기는 추정기는 없다. 능형회귀는 많은 작은 계수가 있다는 믿음, 라쏘는 소수 변수만 중요하다는 믿음을 반영한다. 진실이 다르면 두 방법의 순위도 바뀐다.\par

따라서 알고리즘 이름보다 가정이 중요하다. 예측변수의 정의, 측정오차, 기간 변화, 거래 가능성에 대한 경제적 지식이 모형을 고르는 데 사용된다. 어떤 방법의 최고 성적만 가져와 보편적으로 좋다고 말하면 선택 편향이 생긴다.\par

\subsection{슬라이드 28 · 금융 자료에 맞는 검증}
낮은 신호대잡음비는 작은 $R^2$를 예상하게 하고, 짧은 실효 표본은 계수의 불확실성을 키운다. 시간에 따라 관계가 변하면 임의 분할보다 미래 방향의 검증이 필요하다. 수익률의 공통 충격은 종목 수가 많아도 독립 표본 수가 그만큼 늘지 않음을 뜻한다.\par

이 구조에서 모형의 자유도를 제한하고, 자료의 시점을 엄격히 지키고, 단순한 기준선과 비교하는 것이 과학적 기본값이다. 성과가 실패하면 원인과 예외 기간을 기록해야 한다.\par

\subsection{슬라이드 29 · 편향을 의도적으로 쓰는 처방}
비편향 추정은 참계수를 평균적으로 맞힌다는 장점이 있지만 미래 예측오차를 최소화한다는 보장은 없다. 계수를 0에 가깝게 끌어당기면 약한 신호에서 계수의 요동을 줄일 수 있다. 이 의도적 편향이 \textbf{축소추정}이다.\par

두 처방은 모두 적합과 단순성을 맞바꾼다. 능형회귀는 모든 계수를 연속적으로 줄이고, 라쏘는 일부를 0으로 만든다. 어느 쪽이 맞는지는 실제 변수 효과가 넓게 퍼져 있는지 드문지에 달려 있다.\par

\subsection{슬라이드 30 · 첫 번째 처방으로 가기}
능형회귀는 OLS의 제곱오차에 계수 크기에 대한 벌점을 더한다. 큰 계수가 꼭 나쁘다는 뜻이 아니라, 약한 자료에서 큰 계수를 확신하려면 그만큼 강한 증거가 필요하다는 뜻이다.\par

다음 슬라이드부터 벌점이 왜 불안정한 역행렬을 안정화하는지 유도한다. 절편은 보통 벌점을 주지 않고, 변수는 단위를 맞춰 표준화한다는 실무 조건을 먼저 기억해 두자.\par

\subsection{슬라이드 31 · OLS의 역행렬이 불안정해지는 이유}
$y=Xg+\varepsilon$에서 제곱오차 $(y-Xg)^\top(y-Xg)$를 $g$에 대해 미분하면 $-2X^\top y+2X^\top Xg=0$이다. $X^\top X$가 가역이면 $\hat g=(X^\top X)^{-1}X^\top y$다. 변수들이 서로 비슷하면 $X^\top X$의 어떤 방향의 고유값이 아주 작아진다.\par

추정량의 분산 $\operatorname{Var}(\hat g)=\sigma^2(X^\top X)^{-1}$에서 그 방향의 분산은 $\sigma^2/\lambda_{\min}$이다. 작은 고유값으로 나누니 계수가 폭발한다. $K>N$이면 열들이 독립일 수 없어 해가 유일하지 않으며, 훈련자료를 완벽히 맞히는 여러 해 가운데 미래에 잘 맞는 해를 OLS만으로 고를 수 없다.\par

\subsection{슬라이드 32 · 능형회귀의 해를 직접 유도하기}
능형회귀는 $L(g)=(y-Xg)^\top(y-Xg)+\gamma g^\top g$를 최소화한다. 미분하여 0으로 놓으면 $-2X^\top(y-Xg)+2\gamma g=0$이고, 항을 옮기면 $(X^\top X+\gamma I)g=X^\top y$다. 따라서 다음의 유일한 해를 얻는다.\par

\[
\hat g_{\rm ridge}=(X^\top X+\gamma I)^{-1}X^\top y,\qquad\gamma>0.
\]

$X^\top X$의 고유값 $\lambda_j$에 $\gamma$가 더해지므로 작은 고유값으로 나누는 불안정성이 줄어든다. $\gamma\to0$이면 가역일 때 OLS로, $\gamma\to\infty$이면 표준화된 설명변수의 계수가 0으로 간다. 절편은 별도로 유지되어 과거 평균 예측이 된다.\par

\subsection{슬라이드 33 · 두 변수로 보는 작은 분모}
표준화한 두 변수의 표본 상관이 $\rho$이고 $c_j=X_j^\top y/N$라 하자. 목적함수에 $1/N$을 붙여 미분하면 $(1+\gamma)g_1+\rho g_2=c_1$, $\rho g_1+(1+\gamma)g_2=c_2$다. 두 식을 더하고 빼면 다음이 나온다.\par

\[
g_1+g_2=\frac{c_1+c_2}{1+\rho+\gamma},\qquad
g_1-g_2=\frac{c_1-c_2}{1-\rho+\gamma}.
\]

$\rho=0.99$이고 벌점이 없으면 두 계수의 차이를 $0.01$로 나눈다. 아주 작은 표본 차이가 거대한 상반 계수로 변한다. $\gamma=0.1$이면 분모가 $0.11$이 되어 같은 오차가 훨씬 덜 증폭된다.\par

\subsection{슬라이드 34 · 교점 그림의 해석}
두 예측변수가 거의 같은 방향이면 각각의 계수를 따로 결정할 정보가 부족하다. 회귀의 제곱오차 등고선은 길고 좁은 타원이 되어 작은 데이터 변화에도 최저점이 타원의 긴 축을 따라 크게 움직인다.\par

능형 벌점은 원형 등고선을 더해 원점에서 먼 해에 비용을 부과한다. 결과적으로 같은 두 변수의 합이 가진 신호는 어느 정도 보존하고, 두 변수의 차이에 대한 불필요한 확신을 줄인다. 계수 하나의 해석보다 둘의 합이 더 안정적일 수 있다.\par

\subsection{슬라이드 35 · 두 계수의 숫자 예}
참계수가 둘 다 1인데 표본에서 OLS가 대략 $2.07$과 $-0.08$을 준다면 각각은 크게 틀렸어도 합은 약 2로 맞을 수 있다. 상관이 높은 변수들 사이에서 어느 하나에 신호를 몰아주고 다른 하나를 반대 부호로 보정한 결과다.\par

예측의 목적은 반드시 각 계수의 인과적 진실을 찾는 것이 아니다. 새 자료에서 두 변수가 계속 함께 움직인다면 합의 안정성이 중요하다. 반대로 두 변수가 미래에 따로 움직일 수 있다면 개별 계수의 불확실성도 의사결정에 반영해야 한다.\par

\subsection{슬라이드 36 · 축소와 유효 자유도}
$X^\top X=I$인 직교 조건에서는 $\hat g_{j,\rm ridge}=\hat g_{j,\rm OLS}/(1+\gamma)$다. 모든 계수가 같은 비율로 줄고 정확히 0이 되지는 않는다. 데이터에 맞춘 자유도의 크기를 $d(\gamma)=\operatorname{tr}[X(X^\top X+\gamma I)^{-1}X^\top]$로 정의하면 직교 조건에서는 $K/(1+\gamma)$다.\par

즉 계수는 $K$개지만 각각 표본의 잡음에 완전히 맞춰지지 않는다. 일반적으로 $X^\top X$의 고유값 $\lambda_j$를 쓰면 $d(\gamma)=\sum_j\lambda_j/(\lambda_j+\gamma)$다. 정보가 적은 작은 고유값 방향일수록 더 강하게 축소된다.\par

\subsection{슬라이드 37 · 제약식과 벌점식}
“오차를 최소화하되 $g^\top g\le c$”라는 문제를 생각하자. 라그랑지안은 $(y-Xg)^\top(y-Xg)+\gamma(g^\top g-c)$다. $-\gamma c$는 $g$ 선택에 영향을 주지 않으므로 능형 목적함수와 같다.\par

$c$가 작을수록 원점 가까운 계수만 허용되고 대응하는 $\gamma$는 커진다. 그림에서 OLS의 타원 등고선과 원형 제약 영역이 처음 닿는 점이 능형 해다. 이 기하학은 왜 예측변수에 단위를 맞춰야 하는지도 드러낸다.\par

\subsection{슬라이드 38 · 계수 경로 읽기}
벌점이 아주 크면 모든 기울기 계수는 거의 0이다. 벌점을 줄이면 각 계수가 자료가 가리키는 방향으로 움직인다. \textbf{경로 그림}은 어느 변수가 얼마나 빠르게 커지는지와 상관된 변수들이 서로 값을 나누는지를 보여 준다.\par

의료 자료처럼 신호가 비교적 강한 예제의 경로를 금융에 그대로 대입할 수는 없다. 금융에서는 아주 작은 계수라도 투자 수량을 크게 바꾸면 중요할 수 있고, 표본이 바뀌면 경로가 불안정할 수 있다. 경로는 설명 도구이며 최적 벌점의 판정은 검증자료에 맡긴다.\par

\subsection{슬라이드 39 · 표준화가 필수인 이유}
같은 경제변수를 원 단위로 쓰는지 백만원 단위로 쓰는지에 따라 회귀계수의 수치는 달라진다. 벌점 $\gamma\sum_jg_j^2$는 계수 크기에 직접 적용되므로 단위가 다른 변수에 공평하지 않다.\par

훈련 구간에서 $z_j=(x_j-\bar x_j)/s_j$로 표준화하면 계수는 해당 변수가 한 표준편차 변할 때의 효과가 된다. 검증·시험 자료에도 \textbf{훈련 구간의} $\bar x_j,s_j$를 그대로 적용해야 미래 분포의 정보를 미리 사용하는 누출을 피한다. 절편은 중심화 뒤 별도로 처리한다.\par

\subsection{슬라이드 40 · 편향된 추정기가 왜 이길 수 있는가}
한 계수에 대해 $x^\top x=1$, 참값 $g$, OLS 추정량의 분산 $\sigma^2$라 하자. 능형 추정량은 $\hat g_\gamma=\hat g_{\rm OLS}/(1+\gamma)$다. 기대값은 $g/(1+\gamma)$이므로 편향은 $-\gamma g/(1+\gamma)$, 분산은 $\sigma^2/(1+\gamma)^2$다.\par

\[
\operatorname{MSE}(\gamma)=\frac{\sigma^2+\gamma^2g^2}{(1+\gamma)^2},\qquad
\frac{d\operatorname{MSE}}{d\gamma}
=\frac{2(\gamma g^2-\sigma^2)}{(1+\gamma)^3}.
\]

미분값을 0으로 놓으면 $\gamma^*=\sigma^2/g^2$다. 잡음이 크거나 참계수가 작을수록 강하게 축소한다. 이 최적값은 참 $g$를 모르므로 실제로는 검증자료에서 근사해야 한다.\par

\subsection{슬라이드 41 · 과거 수익률 357개로 다음 달을 예측하기}
한 종목의 과거 119개월 수익률, 그 제곱, 세제곱을 각각 입력하면 357개 특징이 된다. 제곱항은 큰 변동 자체의 효과, 세제곱항은 상승과 하락의 비대칭적 패턴을 표현할 수 있다. 복잡한 기업 정보 없이도 고차원 문제가 된다.\par

매월 종목 간 비교를 한다면 그 달의 횡단면 평균을 빼고 특성을 표준화한다. 이는 시장 전체의 움직임보다 \textbf{어느 종목이 상대적으로 높을지}에 초점을 맞추기 위한 것이다. 변수 정의와 지연 시점은 매매 가능성을 좌우한다.\par

\subsection{슬라이드 42 · 회귀계수에서 읽는 모멘텀과 반전}
최근 몇 달의 양의 계수는 과거에 오른 종목이 더 오르는 모멘텀을 뜻하고, 아주 오래된 과거 수익률의 음의 계수는 장기 반전을 뜻한다. 직전 한 달을 건너뛰는 설계는 호가·유동성에 따른 단기 반전과 중기 모멘텀을 구분하는 데 도움이 된다.\par

다만 계수 곡선은 서로 상관된 시차를 동시에 넣은 \textbf{조건부 효과}다. 한 시차의 계수가 양수라고 그 시차만 사는 전략의 수익률과 동일하지 않다. 전체 예측값을 포트폴리오에 적용해 비용 뒤 성과를 별도로 계산해야 한다.\par

\subsection{슬라이드 43 · 모의실험에서 벌점이 작동하는 방식}
진짜 변수 10개를 넘어서 잡음 변수까지 넣으면 OLS의 표본내 점수는 계속 좋아질 수 있다. 능형 벌점은 잡음 변수의 우연한 계수를 낮춰 새 관측치의 오차를 줄인다. 그렇다고 모든 벌점이 좋은 것은 아니다.\par

과도한 벌점은 진짜 변수까지 거의 0으로 만들어 과거 평균과 다르지 않은 예측을 준다. 따라서 그래프에서는 표본외 곡선이 가장 높은 \textbf{중간 벌점 구간}을 찾는다. 반복 시뮬레이션의 평균 곡선은 한 번의 운 좋은 표본이 아니다.\par

\subsection{슬라이드 44 · 여기까지의 처방}
문제는 변수 수 자체보다 약한 신호에서 자유로운 계수 추정이 큰 분산을 만든다는 점이다. 능형회귀는 계수를 모두 남기면서 특히 정보가 부족한 방향을 강하게 줄인다. 예측에는 유리할 수 있지만 어떤 변수만 실제로 중요한지 알려 주지는 않는다.\par

다음 처방은 많은 변수 중 일부를 정확히 0으로 만드는 라쏘다. 두 방법은 경제적 사전 믿음이 다르므로 성능과 해석을 함께 비교해야 한다.\par

\subsection{슬라이드 45 · 변수 선택의 문제}
재무비율, 가격 신호, 텍스트 지표를 수백 개 넣을 때 모든 계수가 조금씩 비영(非零)인지, 소수만 실질적으로 중요한지 모른다. 변수 선택은 예측 비용을 줄이고 설명을 간결하게 만들 수 있지만, 잘못된 제외는 편향을 키운다.\par

특히 높은 상관을 가진 두 지표가 같은 정보를 담으면 하나만 선택해도 예측은 좋을 수 있다. 하지만 선택된 변수 하나에 유일한 경제적 의미를 부여하면 잘못 해석할 위험이 있다.\par

\subsection{슬라이드 46 · 제곱 벌점에서 절댓값 벌점으로}
능형회귀의 $L_2$ 벌점은 계수를 부드럽게 줄이지만 정확히 0으로 만들지 않는다. 라쏘는 $L_1$ 벌점 $\sum_j|g_j|$을 사용한다. 0에서 꺾이는 모양 때문에 약한 계수는 0에 붙을 수 있다.\par

두 방식 모두 적합도와 복잡성의 균형이다. 차이는 “복잡성”을 계수 제곱합으로 재는지 절댓값 합으로 재는지다. 그 작은 수학적 차이가 변수 선택이라는 질적 차이를 만든다.\par

\subsection{슬라이드 47 · 라쏘 목적함수}
표준화된 변수에서 라쏘는 $\min_g\{N^{-1}\|y-Xg\|^2+\gamma\sum_j|g_j|\}$를 푼다. $\gamma=0$은 OLS와 같고, 커질수록 많은 계수가 0이 된다. 절댓값은 0에서 미분 가능하지 않아 모든 변수에 대한 단순한 역행렬 해는 없다.\par

그러나 볼록 함수라 전역 최저점은 효율적 알고리즘으로 찾을 수 있다. 변수 수가 관측치보다 많아도 유용하며, 해가 하나만 존재하는지 여부는 변수의 상관 구조에 따라 따로 확인해야 한다.\par

\subsection{슬라이드 48 · 연성 임계값의 유도}
직교한 한 변수에서 OLS 계수를 $z$라고 하면 라쏘 문제는 상수항을 빼고 $(g-z)^2+2\lambda|g|$의 최소화와 같다. $g>0$이면 도함수 $2(g-z)+2\lambda=0$이므로 $g=z-\lambda$, 단 $z>\lambda$일 때만 양수다. $g<0$이면 $g=z+\lambda$이고 $z<-\lambda$여야 한다.\par

\[
\hat g_{\rm lasso}=\operatorname{sgn}(z)(|z|-\lambda)_+,
\qquad (u)_+=\max(u,0).
\]

$|z|\le\lambda$이면 0이 최소점이다. 능형회귀는 비율로 줄이고 라쏘는 일정 크기를 빼며, 그 과정에서 약한 계수를 없앤다. 벌점의 정의에 따라 교재의 $\gamma$와 여기의 $\lambda$ 사이에는 2배 상수 차이가 있을 수 있다.\par

\subsection{슬라이드 49 · 마름모의 꼭짓점}
$L_1$ 제약 $|g_1|+|g_2|\le c$는 두 차원에서 마름모이고 꼭짓점이 좌표축에 있다. 제곱오차의 타원 등고선이 제약 영역에 처음 닿는 곳이 꼭짓점이면 한 계수가 정확히 0이다. $L_2$ 제약은 원이므로 축 위의 정확한 접촉이 상대적으로 드물다.\par

기하학은 직관을 주지만 모든 데이터에서 라쏘가 하나의 변수만 남긴다는 뜻은 아니다. 벌점 크기와 변수 상관에 따라 선택 수가 달라지고, 해석을 위해서는 선택의 안정성도 보아야 한다.\par

\subsection{슬라이드 50 · 라쏘 경로의 읽기}
벌점이 매우 크면 기울기 계수가 모두 0이다. 벌점을 줄이면 먼저 자료와 강하게 연결된 변수가 들어오고, 이어 다른 변수가 들어온다. 경로의 각 지점은 하나의 후보 모형이다.\par

변수가 일찍 들어왔다고 인과적으로 중요하다는 뜻은 아니다. 다른 변수와의 상관, 표본의 우연, 스케일에 영향을 받는다. 경로를 비교할 때는 훈련 자료에서만 표준화하고, 성능은 독립 검증에서 평가한다.\par

\subsection{슬라이드 51 · 상관된 쌍둥이 변수}
두 예측변수가 거의 같은 정보를 담으면 라쏘는 하나에 계수를 몰아주고 다른 하나를 버릴 수 있다. 표본이 조금만 바뀌어도 살아남는 변수의 이름이 바뀔 수 있다. 예측은 안정적이어도 변수 선택 결과는 불안정한 것이다.\par

예를 들어 배당/주가와 이익/주가가 같은 가치 신호를 담는다면 하나가 선택되었다고 다른 하나가 경제적으로 무의미하다고 결론 내릴 수 없다. 계수 해석에는 변수 그룹을 함께 살펴야 한다.\par

\subsection{슬라이드 52 · 엘라스틱넷의 두 벌점}
엘라스틱넷은 $\min_g\{N^{-1}\|y-Xg\|^2+\gamma_1\sum_j|g_j|+\gamma_2\sum_jg_j^2\}$를 푼다. $L_1$은 약한 변수를 제거하고 $L_2$는 상관된 변수 사이의 가중치를 나누어 해를 안정화한다.\par

두 벌점을 정해야 하므로 검증 비용이 늘어난다. 하지만 많은 서로 비슷한 재무비율을 사용할 때 선택의 불안정을 낮출 수 있다. ‘둘 다 있으니 무조건 더 낫다’가 아니라, 두 하이퍼파라미터를 선택하는 과정까지 표본외 검증해야 한다.\par

\subsection{슬라이드 53 · 어떤 벌점이 맞는가}
진짜 효과가 넓게 퍼져 있다면 능형회귀의 모든 계수 축소가 유리할 수 있다. 효과가 소수 변수에 집중되면 라쏘의 희소성이 유리하다. 관련 변수들이 군집을 이루면 엘라스틱넷이 실용적일 수 있다.\par

이 판단은 자료가 충분하면 검증 성능으로, 자료가 부족하면 경제적 사전지식으로 보완한다. 여러 방법을 시험한 뒤 최고 성적만 보고하면 모형 선택 자체가 또 하나의 과적합이 된다. 최종 시험자료는 맨 마지막까지 남겨 둔다.\par

\subsection{슬라이드 54 · 벌점 선택과 채점}
벌점 $\gamma$는 회귀계수와 달리 모형을 적합하기 전에 선택하는 \textbf{하이퍼파라미터}다. 같은 훈련자료의 제곱오차만 보면 벌점을 줄일수록 적합이 좋아지므로 항상 0을 선택한다. 미래 성능을 위한 벌점은 보지 않은 자료로 정해야 한다.\par

검증용 기간을 분리하면 훈련 표본이 줄어드는 문제가 있다. 교차검증은 여러 차례 자료를 나누어 각 관측치를 한 번씩 검증에 사용함으로써 제한된 자료를 아낀다.\par

\subsection{슬라이드 55 · 하이퍼파라미터와 계수의 역할}
계수는 벌점을 고정한 뒤 자료에 맞춰 추정한다. 벌점은 얼마나 복잡한 모형을 허용할지 결정한다. 둘을 같은 자료에서 무제한으로 최적화하면 재귀적으로 자유도가 늘어나는 것과 같다.\par

훈련·검증·시험의 역할을 구분한다. 훈련은 계수 추정, 검증은 벌점 선택, 시험은 최종 성능 보고다. 시험 성적을 보며 다시 벌점을 조정하면 시험자료도 더는 독립 표본이 아니다.\par

\subsection{슬라이드 56 · k겹 교차검증을 계산하기}
표본을 $k$개 묶음 $V_1,\ldots,V_k$로 나누고, 각 묶음 $v$를 뺀 자료에서 $\hat g_{-v}(\gamma)$를 적합한다. 뺀 묶음의 오차를 평균한 값이 교차검증 점수다.\par

\[
CV(\gamma)=\frac1k\sum_{v=1}^k\frac1{|V_v|}
\sum_{i\in V_v}(y_i-x_i^\top\hat g_{-v}(\gamma))^2.
\]

이 값을 가장 작게 하는 $\gamma$를 고른다. 묶음 크기가 다르면 전체 관측치로 가중 평균하는 방식도 가능하며, 기준선의 예측을 각 훈련 부분에서 다시 계산해야 공정하다.\par

\subsection{슬라이드 57 · 금융 교차검증의 두 경고}
시간 의존이 있는 수익률을 무작위로 섞으면 미래 관측치가 과거 예측의 훈련에 들어갈 수 있다. 연 단위 블록을 떼거나, 예측 기간이 겹치는 경계 관측치를 제거해 누출을 줄인다. 같은 기업의 여러 달 관측치를 어떻게 묶을지도 목적에 따라 결정한다.\par

교차검증 점수는 벌점 선택에 이미 사용했기 때문에 약간 낙관적이다. 정말 손대지 않은 뒤쪽 기간의 성능이 더 높은 수준의 증거다. 표본내 적합, 검증, 최종 표본외 시험의 정보 순서를 유지한다.\par

\subsection{슬라이드 58 · OLS, 능형, 라쏘의 결과 비교}
같은 과거 수익률 특징 357개를 쓰면 OLS가 표본내 $R^2$는 가장 높아도 검증 점수는 음수가 될 수 있다. 능형과 라쏘는 표본내 적합을 포기하는 대신 검증 점수를 양수로 만들 수 있다. 이것이 편향을 도입한 이유를 경험적으로 보여 준다.\par

다만 표에 적힌 성과는 특정 자료·기간·검증 방식의 결과다. 모형을 실제로 투자에 사용하려면 종목 선정과 매매 시점, 상장폐지 수익률, 거래비용을 포함해 새 기간에서 다시 평가한다. 숫자 하나를 일반 법칙으로 외우지 않는다.\par

\subsection{슬라이드 59 · 예측 점수와 투자 성과가 다른 이유}
$R^2$는 모든 종목의 예측오차 제곱을 벌한다. 포트폴리오는 주로 종목의 상대 순위와 선택 경계 주변의 오류에 민감하다. 예측값의 크기가 다소 틀려도 상위·하위 종목의 순서가 맞으면 롱숏 수익은 괜찮을 수 있다.\par

반대로 $R^2$가 좋아져도 거래가 잦고 유동성이 낮은 종목에 비중을 몰면 실제 성과는 나빠질 수 있다. 그러므로 통계적 목적함수와 경제적 목적함수를 구분한다. 투자비중의 제약과 비용을 반영한 검증은 별도의 설계 문제다.\par

\subsection{슬라이드 60 · 시장 전체 예측으로 돌아오기}
정규화는 시장수익률 예측에서도 변수 선택과 계수의 요동을 줄인다. 그러나 한 나라의 시장 시계열은 독립 역사가 하나뿐이고, 종목 간 횡단면과 달리 같은 달에 수천 개 관측치를 얻지 못한다. 위기와 제도 변화로 관계도 바뀐다.\par

따라서 개별 종목의 상대 수익률 예측에서 개선이 있었다고 시장 타이밍의 큰 개선을 기대할 수는 없다. 강한 기준선과 다양한 시작일에서 견고한지 확인한다.\par

\subsection{슬라이드 61 · 극단적 복잡성의 논쟁}
짧은 자료에 비해 훨씬 많은 무작위 비선형 특징을 만들고 능형회귀를 적용하면, 겉으로는 큰 표본외 성과가 나올 수 있다. 한 해석은 높은 차원의 표현이 놓친 패턴을 잡았다는 것이다. 다른 해석은 특정 설정에서 예측이 최근 수익률의 가중평균으로 수렴해 단순 모멘텀을 재포장한다는 것이다.\par

두 주장을 가르려면 단순한 최근 수익률 규칙과 같은 기간·비용·위험으로 비교하고, 특징 생성과 튜닝의 모든 선택을 포함해 재현해야 한다. 최신 연구의 성공과 반박을 모두 가설로 두는 태도가 적절하다.\par

\subsection{슬라이드 62 · 수익률 예측에서 확립된 규칙}
첫째, 과거 평균이나 0 같은 명확한 기준선을 정한다. 둘째, 데이터의 발표·수정 시점을 지켜 미래 정보가 들어가지 않게 한다. 셋째, 계수의 자유도를 신호대잡음비에 맞춰 제한한다. 넷째, 검증에 사용한 자료와 최종 보고 자료를 분리한다.\par

예측 점수와 투자 수익은 다르므로 어느 결정을 개선하려는지 끝까지 명시한다. 다음 장에서는 시장 전체보다 종목 간 상대 수익률로 대상을 바꾼다. 그때는 이상현상의 경제적 이유와 여러 요인을 동시에 시험하는 문제가 추가된다.\par

\section{04. 종목 간 기대수익률과 요인 논쟁}
시장이 내년에 오를지 예측하는 일과, 같은 달에 어떤 종목이 다른 종목보다 나을지 예측하는 일은 다르다. 후자는 공통 시장 충격을 매수·매도 포지션으로 일부 상쇄할 수 있지만, 작은 주식의 거래비용과 공매도 제약이 중요해진다. 이 장은 종목 특성을 전략으로 만드는 절차에서 출발해 위험 보상·가격오류·우연이라는 세 설명을 평가한다.\par

\subsection{슬라이드 01 · 횡단면의 출발점}
횡단면은 \textbf{한 시점의 여러 종목}을 가리킨다. 시장 시계열 예측은 다음 달 시장 전체의 평균을 묻지만, 횡단면 예측은 같은 달 종목의 상대 순위를 묻는다. 모든 종목이 하락한 달에도 덜 하락한 종목을 사고 더 하락한 종목을 공매도하면 롱숏 차이는 양수일 수 있다.\par

하지만 롱숏 수익은 무위험이 아니다. 공매도 비용, 차입 가능성, 포지션을 다시 맞추는 비용, 시장 급변 시 두 다리의 동시 손실을 함께 고려한다.\par

\subsection{슬라이드 02 · 앞 장에서 건너오는 다리}
시장 전체의 평균수익률은 한 나라의 짧은 역사에서 추정한다. 종목 간 비교에는 매월 여러 기업의 특성이 있어 자료가 풍부해 보인다. 그러나 같은 달 종목은 공통 충격을 공유하므로 표본 행 수가 독립 정보량과 같지는 않다.\par

이번 장의 목표는 종목 특성이 수익률을 예측한다는 주장과, 왜 그런지에 대한 설명을 구분하는 것이다. 다음 단계의 머신러닝은 많은 특성의 함수를 추정하지만 먼저 고전적 실증 결과의 구조를 알아야 한다.\par

\subsection{슬라이드 03 · 요인 동물원의 네 질문}
수백 개 기업 특성 가운데 어느 것이 다른 특성을 통제한 뒤에도 정보를 갖는가. 그 정보를 보상받는 위험으로 설명할 수 있는가. 발표된 전략은 새 기간과 큰 종목에서도 남는가. 여러 변수가 같은 경제적 주제를 반복 측정하는 것은 아닌가.\par

한 특성의 평균 수익률만으로 네 질문에 답할 수 없다. 비교 포트폴리오의 구성, 위험모형, 표본 기간, 거래 가능성을 차례로 명시해야 한다.\par

\subsection{슬라이드 04 · 특성에서 전략으로}
장부가치/시가총액이나 지난 11개월 수익률 같은 특성은 기업별 숫자다. 투자전략은 이 숫자로 종목을 정렬하고 어느 그룹을 사고 팔지 정해야 비로소 정의된다. 같은 특성도 동일가중·시가총액가중, 소형주 포함 여부에 따라 다른 성과를 낸다.\par

또한 특성은 미래 시점에 실제로 알 수 있어야 한다. 연말 재무제표를 그해 12월의 거래에 사용하는 것은 발표 전 정보를 사용한 오류가 될 수 있다.\par

\subsection{슬라이드 05 · 십분위 포트폴리오 정렬}
매월 종목을 특성 $X_{i,t}$ 순서로 정렬해 10개 그룹으로 나누고, 상위 그룹을 매수하고 하위 그룹을 매도한다. 다음 달 수익률 차이 $r_{10,t+1}-r_{1,t+1}$를 계산해 시간 평균을 낸다. 이는 특성의 극단 사이에 수익률 기울기가 있는지 보여 준다.\par

예를 들어 모멘텀은 최근 한 달을 제외한 과거 6\~{}12개월 수익률이 높은 종목을 사는 방식이다. 최근 한 달 제외는 단기 반전과 거래 미시구조의 영향을 줄인다. 실제 전략 성과에는 종목 교체 비용과 공매도 비용을 빼야 한다.\par

\subsection{슬라이드 06 · 파마·맥베스 회귀의 두 단계}
매월 종목들 사이에서 $r_{i,t+1}=a_t+b_tX_{i,t}+e_{i,t+1}$를 회귀하여 그 달의 기울기 $b_t$를 얻는다. 그런 다음 $T$개월 기울기의 평균 $\bar b=T^{-1}\sum_tb_t$가 0과 다른지 검정한다. 여러 특성을 동시에 넣으면 다른 특성을 고정한 조건부 관련성을 볼 수 있다.\par

정렬은 극단 그룹의 차이를 직관적으로 보여 주고 회귀는 여러 특성을 동시에 통제하기 쉽다. 둘 다 수익률과 특성의 관계를 측정할 뿐 원인을 자동으로 밝혀 주지 않는다.\par

\subsection{슬라이드 07 · 높은 수익률과 이상현상의 차이}
고위험 주식이 높은 평균수익률을 주는 것은 위험을 싫어하는 투자자에게 정상적인 보상일 수 있다. \textbf{이상현상}은 특정 기준 자산가격모형이 설명하지 못하는 수익률 차이를 뜻한다. 따라서 기준모형을 바꾸면 이상현상의 유무도 달라질 수 있다.\par

예를 들어 CAPM만 기준으로 할 때의 알파와 시장·가치·규모 요인을 함께 넣은 뒤의 알파는 다르다. “이상”은 자연법칙 위반이 아니라 선택한 모형의 설명 실패라는 의미다.\par

\subsection{슬라이드 08 · 왜 월별 기울기의 표준오차를 쓰는가}
한 달에 3,000개 종목이 있어도 모두 같은 금리·경기 충격을 겪는다. 모든 종목·월을 독립 관측치로 취급한 단순 합동회귀는 표준오차를 지나치게 작게 만들 수 있다. 월별 횡단면 기울기 $b_t$를 구한 뒤 시간에 따라 변하는 $b_t$의 분산으로 평균의 불확실성을 재는 이유다.\par

\[
\widehat{\operatorname{se}}(\bar b)
=\sqrt{\frac{1}{T(T-1)}\sum_{t=1}^T(b_t-\bar b)^2}.
\]

이 식은 월별 기울기가 시간상 독립이라는 단순 조건의 표준오차다. 자기상관이 있으면 Newey–West 같은 장기분산 추정이 필요하다. 독립 표본 수를 종목 수가 아니라 시간의 길이와 연결하는 것이 핵심이다.\par

\subsection{슬라이드 09 · 여섯 범주의 지도}
수익률 예측변수는 모멘텀, 가치, 투자, 수익성, 무형자산, 거래마찰로 묶을 수 있다. 이는 변수를 정리하는 \textbf{분류표}이지 각 범주의 인과이론이 아니다. 같은 범주 안에도 측정 방식과 경제적 채널이 다르다.\par

각 범주를 읽을 때 신호와 매수·매도 방향, 위험 보상 설명, 기대 오류 설명, 재현성과 거래 가능성의 순서로 살펴보면 서로 비교하기 쉽다.\par

\subsection{슬라이드 10 · 여섯 범주의 대표 신호}
모멘텀은 최근 승자, 가치는 낮은 가격/기초가치, 투자는 자산 증가가 낮은 기업, 수익성은 이익률이 높은 기업을 매수하는 방향이다. 무형자산은 연구개발과 조직자본처럼 장부에 잘 보이지 않는 투자, 거래마찰은 규모·유동성·변동성과 연결된다.\par

같은 범주에도 여러 변수와 정렬법이 있다. 한 범주의 일부 전략이 재현된다는 사실을 그 범주의 모든 전략이 유효하다고 해석하면 안 된다. \href{https://doi.org/10.1093/rfs/hhy131}{Hou·Xue·Zhang}은 공통한 구성 규칙으로 대규모 목록을 다시 검정한다.\par

\subsection{슬라이드 11 · 낮은 가격은 두 가지를 말한다}
한 기간 뒤 현금흐름 $CF$의 기대값을 단순화해 쓰면 $P=E[CF]/(1+E[r])$이다. 현재 가격이 낮다는 관찰은 할인율, 즉 요구수익률이 높다는 뜻일 수도 있고 투자자가 미래 현금흐름을 지나치게 비관적으로 봤다는 뜻일 수도 있다.\par

이 식은 위험과 미래 현금흐름의 공분산을 생략한 교육용 단순식이다. 엄밀한 가격식은 확률적 할인인자 $m_{t+1}$을 써서 $P_t=E_t[m_{t+1}CF_{t+1}]$로 나타낸다. 이 구별을 지켜야 회계비율의 예측력을 곧바로 위험 프리미엄이라고 단정하지 않는다.\par

\subsection{슬라이드 12 · 각 범주를 같은 질문으로 읽기}
첫째, 특성은 언제 측정하고 어떤 종목을 매수·매도하는가. 둘째, 왜 수익률 차이가 생기는가. 위험에 대한 보상인지, 투자자의 지나친 낙관·비관이 수정되는 것인지 따진다. 셋째, 원래 표본을 벗어나도 남는가.\par

세 질문은 순서가 있다. 거래를 정의하지 않으면 성과를 측정할 수 없고, 성과가 견고하지 않으면 정교한 설명도 설득력을 잃는다. 설명이 맞더라도 비용 뒤 이익이 없을 수 있다.\par

\subsection{슬라이드 13 · 모멘텀의 위치}
모멘텀은 과거 상승 종목이 단기적으로 더 상승하는 경향이다. 시장 전체가 계속 오르는 것과 다르게 \textbf{같은 달 종목 사이의 상대 승자}를 말한다. 투자자의 정보 반응 지연, 추세 추종 거래, 위험 노출 등이 가능한 설명이다.\par

모멘텀에는 큰 반전 손실이 나타날 수 있어 평균수익률 하나로 위험을 요약하기 어렵다. 다음 세 슬라이드에서 신호·원인·폭락을 분리한다.\par

\subsection{슬라이드 14 · 과거 승자가 계속 이기는가}
매월 형성 시점까지 공개된 가격으로 지난 6\~{}12개월 수익률을 계산하되 직전 한 달을 뺀다. 승자 그룹을 사고 패자 그룹을 판 뒤 다음 달 수익률을 잰다. 실적 서프라이즈 뒤의 가격 지연 반응도 관련 모멘텀 신호다.\par

단기 반전 때문에 직전 한 달을 빼는 설계는 임의적인 숫자 놀이가 아니다. 거래비용·유동성 충격과 중기 정보 반응을 구분하려는 장치다. 대형주에서도 남는지와 매매회전율을 함께 확인해야 한다.\par

\subsection{슬라이드 15 · 모멘텀의 위험과 가격오류 설명}
위험 설명은 승자·패자 포트폴리오가 나쁜 경제 상태에서 손실을 주기 때문에 평상시 높은 평균수익률이 보상이라는 주장이다. 가격오류 설명은 투자자가 새 정보에 천천히 반응해 승자의 좋은 소식이 몇 달에 걸쳐 가격에 반영된다는 주장이다.\par

둘은 폭락 시점과 뉴스 이후 반응에 대해 다른 예측을 한다. 다만 한 사례가 어느 한쪽을 완전히 증명하지는 않는다. 위험 요인이 무엇인지, 투자자가 왜 늦게 반응하는지 각각 검증 가능한 메커니즘을 제시해야 한다.\par

\subsection{슬라이드 16 · 평균 모멘텀과 폭락 위험}
모멘텀의 평균 수익이 양수라도 시장 급락 뒤 강한 반등에서는 과거 패자들이 폭등할 수 있다. 패자를 공매도한 포지션은 이때 큰 손실을 낸다. 평상시 평균과 위기 상태의 손실이 비대칭인 것이다.\par

손실은 단순 표준편차보다 꼬리위험, 마진 요구, 포지션 강제청산에 중요하다. \href{https://doi.org/10.1016/j.jfineco.2015.12.002}{Daniel·Moskowitz}의 분석은 모멘텀 수익과 급격한 반전 위험을 함께 평가하게 한다. 전략의 적정 규모는 평균수익률뿐 아니라 이런 상태에서 버틸 수 있는 자금에 달려 있다.\par

\subsection{슬라이드 17 · 가치와 성장의 구분}
가치주는 주가가 장부가치, 이익, 현금흐름에 비해 낮은 기업이다. “가격이 싸다”는 것은 주당 절대가격이 낮다는 뜻이 아니라 기초가치와의 비율이 낮다는 뜻이다. 성장주는 미래 확장에 대한 기대가 가격에 많이 반영된 기업이다.\par

가치 프리미엄은 가치주를 사고 성장주를 판 상대수익률이다. 두 그룹의 현재 수익성과 미래 성장성은 다를 수 있으므로 단순 비율의 경제적 해석에는 주의가 필요하다.\par

\subsection{슬라이드 18 · 장부가치/시가총액으로 정렬하기}
$B/M$이 높은 종목은 장부상 자기자본에 비해 주가가 낮다. 같은 투자금으로 더 많은 장부가치를 사는 셈이지만, 장부가치의 질이 낮거나 미래 손실이 클 수도 있다. 값이 높다는 사실만으로 저평가를 확정하지 않는다.\par

산업마다 회계처리와 무형자산 비중이 달라 비율이 구조적으로 다르다. 업종 내 비교나 무형자산 보정은 신호를 바꿀 수 있다. 실무에서는 재무제표 발표 지연을 반영하고 상장폐지 종목을 포함해야 성과를 과대평가하지 않는다.\par

\subsection{슬라이드 19 · 가치 프리미엄의 두 해석}
위험 해석은 가치기업의 현금흐름이 경기 악화 때 더 크게 줄어들어 투자자가 높은 기대수익률을 요구한다는 것이다. 가격오류 해석은 시장이 과거 부진을 지나치게 연장해 낮은 가격을 매겼다가 예상을 고치는 과정에서 초과수익이 생긴다는 것이다.\par

둘을 가르려면 나쁜 상태에서의 수익률, 이후 실적발표에 대한 반응, 투자자 기대의 변화가 필요하다. 단순히 $B/M$이 높고 미래 수익률도 높다는 상관 하나로는 충분하지 않다.\par

\subsection{슬라이드 20 · 가치 효과의 기간과 측정}
가치 프리미엄의 크기는 원래 연구 기간, 이후 장기 기간, 대형주만의 결과에 따라 달라질 수 있다. 장부에 계상되지 않은 연구개발·브랜드 자산을 보정하면 현대 기업의 장부가치가 바뀌고 가치 분류도 달라진다.\par

그러나 사후에 가장 높은 성과를 주는 보정법을 고르면 다시 과적합이 생긴다. 회계 원칙으로 보정법을 먼저 정하고 다른 나라·기간에서 확인해야 한다. 무형자산 조정은 단순한 수익률 높이기 기술이 아니라 \textbf{측정오차 개선 가설}이다.\par

\subsection{슬라이드 21 · 투자와 수익성으로 넘어가기}
기업의 실물 투자와 이익 창출 능력은 가격을 결정하는 현금흐름과 할인율 양쪽에 연결된다. 투자율이 높은 기업의 기대수익률이 낮고 수익성이 높은 기업의 기대수익률이 높은 패턴은 처음에는 직관에 어긋나 보인다.\par

하지만 기업이 언제 투자하기를 선택하는지 생각하면 부호를 유도할 수 있다. 이 장의 후반에서 투자 CAPM의 최적화 식으로 다시 확인한다.\par

\subsection{슬라이드 22 · 빠른 자산 증가와 낮은 미래수익률}
총자산이 빠르게 늘어난 기업은 자본비용이 낮아 투자를 많이 할 수 있다. 또는 시장이 확장을 과도하게 낙관한 결과 미래 주식 수익률이 낮아질 수도 있다. 자산 증가율에는 인수합병, 현금 축적, 회계 변화가 섞일 수 있어 지표의 구성을 확인해야 한다.\par

낮은 투자 기업을 사고 높은 투자 기업을 파는 전략의 수익이 양수라면 두 설명 가운데 어느 것이 맞는지 추가 증거가 필요하다. 자금조달 비용, 이후 현금흐름, 발행 시점의 투자자 기대가 도움이 된다.\par

\subsection{슬라이드 23 · 투자율과 할인율의 연결}
새 설비의 기대수익이 할인율보다 높을 때 기업은 투자한다. 다른 조건이 같다면 자본비용이 낮은 기업은 더 많은 프로젝트를 채택할 수 있어 투자율이 높다. 따라서 높은 투자가 낮은 기대주식수익률을 암시할 수 있다.\par

하지만 모든 고투자 기업의 위험이 낮다는 뜻은 아니다. 투자 결정에는 성장 기회, 조정비용, 내부 자금 제약이 함께 들어간다. 투자율은 할인율의 \textbf{간접 신호}이며 다른 조건을 고정해야 해석된다.\par

\subsection{슬라이드 24 · 수익성 높은 회사의 높은 수익률}
이익/자산이나 자기자본이익률이 높은 기업이 미래에도 높은 수익률을 주는 패턴은 “좋은 기업은 이미 비쌀 텐데 왜 더 오르는가”라는 질문을 만든다. 가격이 비싸도 미래 현금흐름이 시장의 예측보다 더 좋아지면 추가 수익이 가능하다.\par

또는 높은 기대 수익성이 같은 투자율 아래에서 높은 할인율과 함께 나타날 수 있다. 수익성은 회계상 일회성 이익과 산업별 자본집약도에 영향을 받으므로 지속 가능한 영업 성과로 측정해야 한다.\par

\subsection{슬라이드 25 · 수익성 효과의 항등식과 설명}
기업가치를 미래 현금흐름의 현재가치로 쓰면, 현재 가격과 예상 투자 규모가 주어졌을 때 기대 이익이 클수록 기대수익률도 높아져야 한다. 이는 특정 위험의 존재를 증명하는 이론이 아니라 가격·현금흐름·수익률 사이의 회계적 관계다.\par

위험 설명은 높은 수익성이 일정 투자와 함께 나타날 때 높은 할인율을 반영한다는 것이다. 가격오류 설명은 투자자가 우량 기업의 지속성을 충분히 가격에 반영하지 못한다는 것이다. 대형 유동주에서도 효과가 남는지, 이후 이익 서프라이즈가 있는지 비교한다.\par

\subsection{슬라이드 26 · 두 자산가격 요인모형}
파마·프렌치 5요인모형은 시장, 규모, 가치에 수익성 RMW와 투자 CMA를 더한다. 다른 q 요인모형도 수익성과 투자를 중심에 둔다. 요인을 추가하면 원래 이상현상의 알파가 줄어들 수 있지만, 이것만으로 요인이 경제적으로 올바른 위험이라는 뜻은 아니다.\par

요인모형은 넓은 종목 집합의 평균수익률을 소수 공통 변수로 설명하는 도구다. 모델 간 경쟁은 어떤 포트폴리오를 기준으로 평가하고, 새 자료에서도 가격오차가 작은지로 이루어져야 한다.\par

\subsection{슬라이드 27 · 무형자산의 회계 문제}
기계는 자산으로 올리고 여러 해에 걸쳐 감가상각하는 반면 연구개발·조직 구축 비용은 즉시 비용 처리되는 경우가 많다. 경제적으로는 둘 다 미래에 수익을 줄 수 있는 투자여도 회계 장부에는 다르게 남는다.\par

따라서 장부가치, 수익성, 투자율의 비교가 무형자산 비중 높은 기업에 불리하게 왜곡될 수 있다. 단, 모든 연구개발 지출이 성공적 자산으로 남는다는 가정도 과하다.\par

\subsection{슬라이드 28 · 장부 밖 투자와 비율의 왜곡}
두 기업이 각각 100을 미래 생산능력에 투입한다고 하자. 설비 구매는 자산에 남고 당기 감가상각 10만 비용일 수 있지만, 연구개발 100은 당기 비용으로 처리될 수 있다. 그러면 경제적 투자 규모가 비슷해도 보고 이익과 장부자본이 크게 다르다.\par

이 차이는 $B/M$과 ROE의 분자·분모를 동시에 바꾼다. 무형자산을 자본화하려면 내용연수와 상각률을 가정해야 하므로 새 측정오차도 생긴다. 회계 보정의 경제적 근거와 민감도 분석을 함께 제시해야 한다.\par

\subsection{슬라이드 29 · 무형자산 프리미엄의 두 설명}
핵심 인력과 기술 지식은 경기 충격 때 가치가 불안정할 수 있어 위험 보상을 요구할 수 있다. 다른 한편 투자자가 장부에 드러나지 않는 역량을 늦게 인식하면 가격오류가 생길 수 있다. 연구개발의 수익은 성공과 실패가 크게 갈려 평균만으로 판단하기 어렵다.\par

고용 이동, 특허, 광고, 영업레버리지 같은 대용변수는 서로 다른 무형자본을 측정한다. 여러 지표를 하나의 “무형자산 효과”로 묶을 때는 각각의 측정 대상이 무엇인지 확인한다.\par

\subsection{슬라이드 30 · 낮은 범주 재현율을 해석하기}
무형자산 범주의 많은 신호가 공통 기준에서 재현되지 않는다고 해서 무형자본 자체가 수익률과 무관하다고 결론 낼 수 없다. 범주에 분석가 전망, 이익의 질, 계절성 같은 변수가 함께 포함되기 때문이다.\par

연구개발을 매출로 나누는지 시가총액으로 나누는지에 따라 신호의 의미도 달라진다. 전자는 활동의 강도, 후자는 시장이 붙인 가격 대비 지출을 뜻한다. 범주의 재현율은 \textbf{분류표의 평균}이지 개별 경제 기제의 검정이 아니다.\par

\subsection{슬라이드 31 · 거래마찰은 가격과 비용을 함께 바꾼다}
규모가 작고 거래가 드문 종목의 높은 평균수익률은 위험 프리미엄일 수도 있지만, 실제 투자자가 그 가격에 거래하기 어려운 데 대한 보상일 수도 있다. 호가 스프레드와 가격 충격을 빼면 종이 위의 이익이 줄어든다.\par

따라서 실증 전략의 수익률을 읽을 때 수익을 만든 종목의 시가총액 비중과 매매 가능한 규모를 함께 본다. 동일가중 포트폴리오는 작은 기업을 경제 전체에서 차지하는 비중보다 크게 반영한다.\par

\subsection{슬라이드 32 · 소형주와 비유동성}
소형주는 분석가의 관심이 적고 정보가 가격에 늦게 반영될 수 있다. 동시에 경기 악화 때 자금조달이 더 어렵고 거래를 청산하기도 힘들어 투자자가 보상을 요구할 수 있다. 두 설명은 함께 성립할 수 있다.\par

거래가 드문 주식의 마지막 체결가격은 현재 시장가치보다 오래된 값일 수 있다. 관측 수익률과 실제 가능한 매도가격의 차이를 무시하면 작은 종목 중심 전략의 성과가 과장된다.\par

\subsection{슬라이드 33 · 고유변동성과 단기 반전}
시장요인으로 설명되지 않는 개별 종목의 변동성을 고유변동성이라 한다. 전통적 분산투자 이론에서는 잘 분산할 수 있는 위험에 높은 보상이 필요하지 않지만, 실제 자료에서는 고유변동성이 높은 종목의 낮은 평균수익이 관찰되기도 한다.\par

직전 한 달 급등·급락 뒤 반전은 투자자 과잉반응 또는 유동성 공급의 대가일 수 있다. 호가 반등처럼 미시구조에서 나온 가짜 패턴도 가능하다. 비용과 대형주 표본에서 남는지를 확인해야 한다.\par

\subsection{슬라이드 34 · 거래마찰 범주의 재현}
소형·비유동주를 큰 비중으로 담는 전략은 원래 표본에서는 강한 차이를 보일 수 있지만 시가총액가중이나 대형주 기준으로 바꾸면 약해진다. 이는 통계적 존재와 실제 운용 가능성이 다른 질문임을 보여 준다.\par

높은 평균 수익률이 거래비용의 보상이라면 비용을 빼기 전 수익은 ‘공짜 돈’이 아니다. 매매 규모가 커질수록 자신의 거래가 가격을 움직인다는 시장충격도 반영해야 한다.\par

\subsection{슬라이드 35 · 여섯 범주를 한꺼번에 비교하기}
모멘텀은 평균 성과와 폭락 위험, 가치는 할인율과 예상오류의 구분, 투자는 기업의 최적 선택, 수익성은 현금흐름과 가격의 관계, 무형자산은 측정, 거래마찰은 구현 비용이 핵심이다. 서로 다른 범주에 한 가지 설명을 억지로 적용하면 중요한 차이가 사라진다.\par

재현율은 해당 범주의 여러 지표 중 정해진 통계 기준을 넘은 비중이다. 이것은 경제적 원인이나 모든 종목에서의 거래가능 수익률을 뜻하지 않는다. 다음 절에서 위험·가격오류·우연을 구분하는 검증을 살핀다.\par

\subsection{슬라이드 36 · 위험인가 가격오류인가}
같은 예측변수가 미래 수익률과 상관되어도 경제적 해석은 갈린다. 위험 프리미엄이라면 나쁜 상태에서 손실을 주는 보상이고, 가격오류라면 투자자의 기대가 체계적으로 잘못되었다가 수정되는 이익이다.\par

이 구분은 투자자에게도 중요하다. 위험 보상은 시장이 효율적이어도 남을 수 있지만 위기 손실을 감수해야 한다. 가격오류는 경쟁과 자금 유입으로 줄어들 수 있으며, 그 속도는 차익거래 제약에 달려 있다.\par

\subsection{슬라이드 37 · 두 설명의 서로 다른 예측}
위험 설명은 전략의 수익이 소비가 낮거나 유동성이 부족한 상태에서 나빠진다고 예측한다. 가격오류 설명은 잘못된 기대가 드러나는 실적발표 같은 시점에 수익이 집중되고, 공개·발표 후에는 효과가 감소할 수 있다고 예측한다.\par

하지만 발표 후 수익 감소만으로 가격오류를 확정할 수는 없다. 위험의 가격이 바뀌거나 기업의 구성이 변했을 수 있다. 하나의 검정이 아니라 여러 시점·상태·집단에서 예측을 맞춰 보아야 한다.\par

\subsection{슬라이드 38 · 지금까지의 가격 항등식}
현재의 낮은 가격은 미래 현금흐름에 대한 낮은 예상 또는 높은 할인율과 양립한다. 종목 특성이 높은 미래 수익률을 예측한다면 먼저 \textbf{어느 경로를 통해 가격과 연결되는지} 물어야 한다.\par

이전 절의 여섯 범주는 같은 항등식을 서로 다른 변수로 시험한다. 모멘텀은 정보 반응의 시간, 무형자산은 현금흐름 측정, 투자는 기업의 자본비용 선택을 특히 강조한다.\par

\subsection{슬라이드 39 · 네 가지 구별 질문}
발표 뒤 수익이 줄어드는가, 이익이 매수 다리와 공매도 다리 중 어디에 있는가, 큰 유동주에도 남는가, 기업 특성보다 위험요인과의 공분산이 중요한가를 묻는다. 각각은 가격오류와 위험 보상에 대해 다른 힌트를 준다.\par

어느 질문도 혼자서 결정타가 되지는 않는다. 공매도 다리의 수익은 공매도 제약 때문에 과대평가가 오래 남는다는 해석과 맞지만, 그 종목의 높은 차입 비용이 이익을 지울 수도 있다.\par

\subsection{슬라이드 40 · 판정표의 증거를 읽는 법}
발견 이후 평균 롱숏 수익률이 감소했다는 연구, 투자자 심리가 높은 때 공매도 다리의 성과가 크다는 연구, 대형주에서도 모멘텀이 남는다는 연구는 서로 다른 증거를 제공한다. 기간과 전략 구성이 달라 숫자를 한 표에 놓고 직접 대소 비교하면 안 된다.\par

특성과 위험요인 노출을 분리한 실험도 표본에 민감하다. 표에서 유의한 결과가 있더라도 경제적으로 예상되는 방향, 다른 표본에서의 재현, 매매 가능성을 차례로 본다. \href{https://doi.org/10.1111/jofi.12365}{McLean·Pontiff}는 발표 전후의 변화를 체계적으로 살핀다.\par

\subsection{슬라이드 41 · 예측자와 설명자의 목표}
예측모형은 장부가치/시가총액이 다음 수익률을 맞힌다면 위험 때문인지 가격오류 때문인지 몰라도 사용할 수 있다. 그러나 투자자 행동이 바뀐 뒤에도 신호가 남을지 예측하려면 원인을 이해해야 한다.\par

또한 공공정책이나 규제 효과는 과거 상관만으로 평가하기 어렵다. 예측에 유용한 변수와 시장의 구조적 메커니즘을 구분하는 것이 수익률 연구의 두 과제다.\par

\subsection{슬라이드 42 · 투자 CAPM의 최적화 유도}
기업이 현재 자산 $A$에 추가 투자 $I$를 하고, 투자 한 단위의 조정비용이 투자율 $I/A$에 따라 올라간다고 하자. 현재 총비용을 $I+(a/2)(I/A)^2A$로 쓰면 $I$에 대한 한계비용은 $1+aI/A$다. 한계 투자에서 얻는 미래 이익 $\Pi_1$의 현재가치를 $E[M_1\Pi_1]$라 하면 최적화 조건은 다음과 같다.\par

\[
1+a\frac{I}{A}=E[M_1\Pi_1].
\]

위험과 할인인자를 단순화하여 한 기간 요구 총수익률 $1+E[r]$로 현재가치를 표현하면 $E[M_1\Pi_1]\approx E[\Pi_1]/(1+E[r])$다. 따라서 $1+E[r]\approx E[\Pi_1]/(1+aI/A)$다. 예상 수익성이 같다면 높은 투자는 낮은 요구수익률과, 투자율이 같다면 높은 예상 수익성은 높은 요구수익률과 연결된다. 이 최적화는 부호를 설명하지만 \textbf{왜 할인율이 기업마다 다른지}는 별도의 위험·시장 이론이 필요하다.\par

\subsection{슬라이드 43 · 특성과 공분산을 떼어 놓기}
가치주가 높은 수익을 얻는 이유가 가치요인 HML에 대한 노출이라면, 장부가치/시가총액을 고정하고 HML 노출만 높은 종목은 더 높은 수익을 주어야 한다. 반대로 특성 자체가 중요하다면 노출 차이를 만들어도 수익 차이가 작을 수 있다.\par

실험은 규모와 가치 특성이 비슷한 종목 묶음 안에서 위험요인 노출로 다시 정렬한다. 이런 조건부 비교는 단순 상관보다 강하지만, 노출 추정의 잡음과 표본 기간에 민감하다. 서로 다른 장기 표본이 반대 결과를 내면 한 연구로 판정이 끝나지 않는다.\par

\subsection{슬라이드 44 · 우연이라는 세 번째 설명}
높은 수익률과 그럴듯한 이야기가 함께 있어도 많은 후보 가운데 운 좋게 발견된 결과일 수 있다. 논문에 실린 전략만 보면 실패한 시도는 보이지 않는다. 이 \textbf{선택된 표본의 문제}가 요인 동물원의 핵심이다.\par

가격오류와 위험 프리미엄을 논하기 전에 신호가 우연인지 검토해야 한다. 연구의 전체 탐색 과정, 관련 변수의 수, 독립 기간의 성적을 가능한 한 기록한다.\par

\subsection{슬라이드 45 · t값 2의 의미가 바뀌는 순간}
한 개의 사전 지정된 가설을 정규근사로 검정하면 $|t|>1.96$은 양측 5\% 수준이다. 하지만 효과가 0인 전략을 300개 시도하면 평균 약 15개가 그 기준을 넘을 수 있다. 가장 높은 값만 발표하면 독자는 300번의 탐색을 보지 못한다.\par

또한 $t$값은 평균/표준오차이므로 경제적 크기와 다르다. 거래비용을 넘는 월수익률인지, 위기 손실을 감수할 만한지, 기간을 늘려도 같은 부호인지 따로 봐야 한다.\par

\subsection{슬라이드 46 · 잡음 전략 300개의 모의실험}
참 기대수익률이 모두 0인 300개 전략을 독립적으로 만들고 각 전략의 평균에 대한 $t$값을 계산하면 히스토그램의 일부 막대가 우연히 2를 넘는다. 이때 발견된 ‘이상현상’은 모의실험 설계상 모두 거짓이다.\par

독립 가정이 현실에서는 성립하지 않아 15라는 기대 숫자가 그대로 맞지 않을 수 있다. 하지만 여러 번 시도한 뒤 큰 값만 고르면 거짓 발견이 생긴다는 원리는 변하지 않는다. 임계값을 높이거나 탐색되지 않은 자료로 다시 검정한다.\par

\subsection{슬라이드 47 · 진짜 전략 10개를 섞었을 때}
잡음 290개 사이에 참 신호가 있는 전략 10개를 넣으면 낮은 임계값은 진짜도 많이 잡지만 거짓도 많이 포함한다. 높은 임계값은 거짓을 줄이는 대신 약한 참 신호를 놓친다. 이것이 검정력과 거짓 발견의 상충관계다.\par

정책 목적에 따라 기준이 달라진다. 수백 후보에서 연구할 만한 주제를 고르는 단계라면 거짓 발견 비율을 통제할 수 있고, 고객 자금을 큰 규모로 투입한다면 더 엄격한 독립 확인이 필요하다.\par

\subsection{슬라이드 48 · 왜 새 요인에는 더 높은 기준이 필요한가}
발표된 요인이 수백 개라면 각 논문의 5\% 유의성 기준을 그대로 유지하기 어렵다. Harvey·Liu·Zhu는 다중검정을 고려해 새 요인의 발견 기준을 대략 $t>3$으로 높일 것을 제안했다.\par

숫자 3은 모든 연구에 똑같이 적용되는 자연상수가 아니다. 실제로 몇 번 시도했는지, 전략들 사이의 상관, 사전 가설의 강도에 따라 적절한 기준이 달라진다. \href{https://doi.org/10.1093/rfs/hhv059}{연구 원문}의 수치는 요인 목록이 커질 때 기존 판정 규칙을 재고해야 함을 보여 준다.\par

\subsection{슬라이드 49 · 다중검정 규칙의 단계별 계산}
$M$개 검정에서 하나라도 거짓 양성이 생길 확률을 5\% 이하로 두고 싶다면 합집합의 확률은 개별 확률의 합 이하라는 부등식을 쓴다. 각 검정을 $\alpha/M$에서 실시하면 전체 거짓 양성 확률은 최대 $M(\alpha/M)=\alpha$다. 이것이 본페로니 기준이다.\par

\[
P\!\left(\bigcup_{j=1}^{M}A_j\right)
\le\sum_{j=1}^{M}P(A_j),\qquad
p_j\le\frac{\alpha}{M}.
\]

Holm은 $p$값을 작은 순서로 정렬해 $p_{(k)}\le\alpha/(M-k+1)$를 차례로 확인한다. Benjamini–Hochberg는 가족 전체에서 단 한 개의 오류도 허용하지 않는 대신, 발견된 것 가운데 거짓의 기대 비율을 통제한다. 독립 또는 특정 의존 조건 아래 최대 $k$에 대해 $p_{(k)}\le k\alpha/M$을 적용한다. 어떤 오류를 통제하려는지 먼저 정해야 방법을 고를 수 있다.\par

\subsection{슬라이드 50 · 재현 논쟁의 질문}
한 연구는 대부분의 이상현상이 공통한 엄격한 조건에서 실패한다고 보고하고, 다른 연구는 상당수가 다양한 나라와 기간에서 남는다고 보고한다. 모순처럼 보이지만 “무엇을 같은 전략으로 볼지”와 “재현을 무엇으로 정의할지”가 다르다.\par

원 논문의 숫자 하나만 비교하는 대신 포트폴리오 가중, 소형주 처리, 알파의 기준모형, 변수 목록의 선택, 통계 기준을 나누어 보아야 한다.\par

\subsection{슬라이드 51 · 많은 이상현상이 실패한다는 결과}
Hou·Xue·Zhang은 여러 전략을 동일한 규칙으로 재구성하고 NYSE 종목의 크기를 기준으로 그룹을 나누며 시가총액 가중 수익률을 쓴다. 이 방식에서는 작은 기업이 종목 수에 비해 시장가치에서 차지하는 비중이 작으므로 그 효과가 전체 전략을 지배하기 어렵다.\par

상당수 신호가 통계적 기준을 넘지 못한다는 결과는 “어떤 수익 차이도 존재하지 않는다”가 아니다. 원래 동일가중·초소형주 중심의 차이는 있을 수 있으나 큰 자금을 운용하는 투자자가 얻을 수 있는 성과와는 다르다는 뜻이다.\par

\subsection{슬라이드 52 · 많은 요인이 재현된다는 결과}
Jensen·Kelly·Pedersen은 요인을 더 적은 주제로 묶고 여러 나라의 자료까지 살펴 높은 재현율을 보고한다. 한 주제의 여러 변형은 완전히 독립된 20개 발견이 아니라 같은 잠재 효과에 대한 여러 측정일 수 있다.\par

이 관점에서는 관련 요인의 정보를 함께 쓰는 계층모형이 유용하다. 그러나 같은 잘못된 데이터 처리에서 나온 변형들이 서로를 강화할 위험도 있다. 발표 뒤와 다른 나라에서의 성과가 중요한 외부 확인이다. \href{https://doi.org/10.1111/jofi.13249}{연구 원문}을 통해 비교 기준을 확인할 수 있다.\par

\subsection{슬라이드 53 · 왜 판정이 갈리는가}
원수익률 차이를 볼지 CAPM 알파를 볼지, 십분위와 삼분위 중 무엇을 쓸지, 소형주를 어떻게 제한할지, 전체 목록인지 이미 선별한 목록인지가 다르다. 또한 단일 요인마다 고정 임계값을 적용할지 주제별 증거를 함께 사용할지가 다르다.\par

따라서 두 재현율을 한 식에 넣어 평균내면 안 된다. 어떤 투자자에게 중요한 질문인지 정한 뒤 그에 맞는 구성 규칙과 손실함수를 택한다. 초소형주를 거래할 수 없는 기관투자자와 소규모 연구 포트폴리오의 답이 다를 수 있다.\par

\subsection{슬라이드 54 · 한 요인의 베이지안 축소를 유도하기}
진짜 알파가 $\alpha\sim N(0,\tau^2)$이고 관측 추정치가 $\hat\alpha\mid\alpha\sim N(\alpha,\sigma^2/T)$라고 하자. 정규분포의 사전밀도와 우도를 곱해 $\alpha$의 제곱항을 완전제곱하면 사후 평균은 정밀도 가중평균이 된다.\par

\[
E[\alpha\mid\hat\alpha]
=\frac{T/\sigma^2}{1/\tau^2+T/\sigma^2}\hat\alpha
=\frac{\tau^2}{\tau^2+\sigma^2/T}\hat\alpha.
\]

자료가 길어 $T$가 크면 계수가 1에 가까워 관측치를 믿는다. 사전분포가 0 근처에 강하게 모이거나 추정 잡음이 크면 0에 더 가깝게 축소한다. 이는 능형회귀와 같은 \textbf{추정오차를 줄이려는 논리}가 요인 알파에도 적용된 것이다.\par

\subsection{슬라이드 55 · 주제별 정보를 함께 쓰기}
요인 $i$의 참 알파를 해당 주제의 공통효과 $c_{j(i)}$와 요인 고유의 편차 $w_i$의 합으로 둔다. $c_j\sim N(0,\tau_c^2)$, $w_i\sim N(0,\tau_w^2)$라는 가정은 주제 안의 요인들이 정보를 공유하도록 한다.\par

자료가 주제 평균을 강하게 지지하면 개별 요인은 같은 방향으로 일부 힘을 얻는다. 반대로 주제 전체가 약하면 우연히 큰 요인도 0 쪽으로 당겨진다. 이런 ‘정보 빌리기’는 유용하지만 주제 분류 자체가 틀리면 오히려 편향을 퍼뜨릴 수 있다.\par

\subsection{슬라이드 56 · 사후 판정과 82\%의 해석}
사후 평균을 0으로 축소하면 유의한 요인 수가 줄어들 수 있지만, 관련 요인의 정보를 함께 쓰면 사후 불확실성도 낮아진다. 두 효과가 상쇄되면 단순 알파 검정과 계층모형의 재현율이 비슷할 수 있다.\par

높은 재현율을 베이지안 방법만의 성과로 해석하지 말아야 한다. 선별된 요인 목록, CAPM 알파라는 목표, 고정된 높은 $t$ 기준의 부재가 함께 작용한다. \href{https://doi.org/10.1111/jofi.13249}{JKP의 연구}를 읽을 때 이 선택을 따로 비교한다.\par

\subsection{슬라이드 57 · 발표 뒤 수익의 감소}
전략의 원 논문 표본, 표본 종료와 발표 사이, 발표 이후를 분리하면 감소가 언제 시작했는지 볼 수 있다. 표본 밖에서 이미 줄었다면 처음 추정값에 선택 편향이 있었을 수 있다. 발표 후 추가 감소는 투자자의 학습과 거래 경쟁에 부합한다.\par

하지만 평균 변화에 시대별 위험 프리미엄과 거래 가능 종목의 구성이 섞일 수 있다. 발표 후 성과가 양수로 남더라도 비용 뒤 이익이 남는지와 나쁜 상태의 위험은 별도 문제다.\par

\subsection{슬라이드 58 · 새 요인의 올바른 검정}
새 요인을 파마·프렌치 3요인만 통제한 뒤 유의하다고 선언하면 이전에 발견된 많은 요인과 겹치는 정보를 ‘새로운’ 것으로 착각할 수 있다. 전체 기존 요인을 모두 통제하면 변수 수가 너무 많고 상관이 높아 OLS가 불안정해진다.\par

정규화가 필요하지만, 관심 대상은 단순 예측점수가 아니라 새 요인의 \textbf{독립적인 위험가격}이다. 통제변수를 선택하는 방법이 그 계수의 편향을 만들지 않도록 설계해야 한다.\par

\subsection{슬라이드 59 · 선택된 통제만 넣을 때의 편향}
기존 요인 $Z$가 평균수익률 $y$에는 약하게 연결되더라도 새 요인 $D$와 강하게 상관될 수 있다. $Z$를 빠뜨리고 $y$를 $D$에만 회귀하면 $D$가 $Z$의 효과를 대신 가져간다. 이는 미시계량의 \textbf{누락변수 편향}과 같은 구조다.\par

한 번의 라쏘로 $y$를 잘 예측하는 통제만 고르면 이런 $Z$를 놓칠 수 있다. 새 요인의 독립성을 판단하려면 평균수익률을 설명하는 변수뿐 아니라 새 요인과 연결된 변수도 찾아야 한다.\par

\subsection{슬라이드 60 · 새 요인 15개의 비교}
좁은 3요인 기준에서 유의하던 새 요인 중 일부는 기존 요인 전체를 적절히 통제하면 약해진다. 반대로 기존 요인과 구별되는 수익성 신호 등은 남을 수 있다. “유의하지 않음”은 효과가 정확히 0이라는 증명이 아니라 그 검정의 불확실성이 크다는 의미다.\par

표의 서로 다른 $t$값은 같은 수익률을 서로 다른 통제집합 아래에서 측정한 것이다. 요인의 이름보다 어떤 위험을 추가로 설명하는지, 이전 요인과 얼마나 중복되는지 살핀다.\par

\subsection{슬라이드 61 · 이중 선택의 세 단계}
평균수익률 벡터를 $\bar r$, 기존 요인 노출 행렬을 $C_h$, 새 요인 노출을 $C_g$라 하자. 첫 라쏘는 $\bar r$를 $C_h$로 예측해 가격 결정에 중요한 통제 $I_1$을 고른다. 둘째 라쏘는 $C_g$를 $C_h$로 예측해 새 요인과 상관된 통제 $I_2$를 고른다.\par

\[
I_1=\operatorname{supp}\arg\min_b\{\|\bar r-C_hb\|^2+\lambda_1\|b\|_1\},\qquad
I_2=\operatorname{supp}\arg\min_c\{\|C_g-C_hc\|^2+\lambda_2\|c\|_1\}.
\]

마지막에는 $C_g$와 $I_1\cup I_2$의 통제를 넣고 벌점 없는 회귀로 새 요인의 계수를 추정한다. 둘째 선택이 빠지면 결과에는 약하지만 새 요인과 강하게 연결된 변수가 누락될 수 있다. \href{https://doi.org/10.1111/jofi.12883}{Feng·Giglio·Xiu}는 이중 선택을 요인 검정에 적용한다.\par

\subsection{슬라이드 62 · 다음 문제의 문턱}
수백 기업 특성의 어느 조합이 다음 달 수익률을 가장 잘 예측하는가라는 질문이 남는다. 고전적 정렬과 회귀는 개별 신호의 경제적 의미를 보여 주지만 비선형 상호작용이 많으면 하나씩 시험하기 어렵다.\par

그렇다고 더 복잡한 알고리즘이 자동으로 답이 되는 것은 아니다. 앞 장의 시간 순서 검증과 이번 장의 다중검정·거래가능성 기준을 그대로 적용해야 한다.\par

\subsection{슬라이드 63 · 요인 논쟁을 예측 문제로 다시 쓰기}
수학적으로 목표는 $E[r_{i,t+1}\mid X_{i,t}]$라는 조건부 평균 함수다. $X_{i,t}$에는 기업 규모, 장부비율, 모멘텀, 수익성 등 시점 $t$에 알려진 특성이 들어간다. 선형 요인모형은 이 함수에 강한 구조 제약을 두는 한 방식이다.\par

유연한 머신러닝은 상호작용과 비선형성을 허용하지만 예측 함수의 경제적 설명과 과적합 통제가 더 중요해진다. 조건부 평균을 맞히는 것과 위험 프리미엄의 구조를 밝히는 것은 서로 다른 목표다.\par

\subsection{슬라이드 64 · 이미 시도한 정규화의 교훈}
앞 장에서는 과거 수익률의 여러 시차와 제곱·세제곱을 넣어 능형회귀와 라쏘를 비교했다. OLS의 표본내 우수성이 검증자료에서는 유지되지 않았고, 계수 축소가 개선을 가져왔다. 이 결과는 새로운 특성 집합에도 같은 평가 규칙이 필요함을 보여 준다.\par

단, 이전의 입력은 과거 가격뿐이었다. 재무제표와 기업 특성이 추가되면 정보의 발표 시점, 산업별 측정 차이, 상장폐지 자료를 더 세심하게 다뤄야 한다.\par

\subsection{슬라이드 65 · 횡단면 연구의 결론}
기업 특성에서 롱숏 수익을 만드는 방법과 그 차이를 이상현상이라 부르는 조건을 구분했다. 모멘텀·가치·투자·수익성·무형자산·거래마찰은 평균수익률 패턴을 조직하는 언어이며, 위험 보상·가격오류·우연을 각각 시험해야 한다.\par

재현 논쟁은 자료와 질문의 정의가 결론을 바꾼다는 사실을 보여 준다. 새 요인은 기존의 많은 요인과 비교하고, 독립 기간과 거래비용 뒤에도 의미가 있는지 확인한다. 예측력의 발견은 엄격한 검증과 경제적 해석이 함께 갈 때 비로소 설득력을 갖는다.\par

\end{document}
